% 使用 XeLaTeX 编译；默认生成含答案版。
% 若需生成练习版，请在下一行取消注释后编译：
% \def\PracticeEdition{1}
\documentclass[UTF8,openany,oneside,zihao=-4,fontset=windows]{ctexbook}
\setmainfont{DejaVu Serif}
\setmonofont{DejaVu Sans Mono}
\usepackage[a4paper,top=2.3cm,bottom=2.3cm,left=2.35cm,right=2.35cm,headheight=15pt]{geometry}
\usepackage{amsmath,amssymb,bm}
\usepackage{array,booktabs,longtable,tabularx}
\usepackage{enumitem}
\usepackage{xcolor}
\usepackage{etoolbox}
\usepackage[most]{tcolorbox}
\usepackage{fancyhdr}
\usepackage{needspace}
\usepackage{fvextra}
\usepackage{hyperref}
\usepackage{bookmark}
\usepackage{microtype}

\definecolor{MVBlue}{HTML}{1F4E79}
\definecolor{MVLight}{HTML}{EDF4FA}
\definecolor{MVGold}{HTML}{C58A18}
\definecolor{MVGray}{HTML}{58636D}
\definecolor{MVCode}{HTML}{F5F7F9}

\hypersetup{colorlinks=true,linkcolor=MVBlue,urlcolor=MVBlue,pdfauthor={NKU AI Study},pdftitle={机器视觉技术练习题合集},pdfsubject={机器视觉技术复习题}}
\setlength{\parindent}{2em}
\setlength{\parskip}{0.35em}
\linespread{1.18}
\setlist{nosep,leftmargin=2.2em}

\newif\ifshowsolutions
\ifdefined\PracticeEdition
  \showsolutionsfalse
\else
  \showsolutionstrue
\fi

\ctexset{
  chapter={format=\huge\bfseries\color{MVBlue},number=\chinese{chapter},name={第,编},beforeskip=0pt,afterskip=1.2em},
  section={format=\Large\bfseries\color{MVBlue},beforeskip=1.3em,afterskip=0.65em},
  subsection={format=\large\bfseries\color{MVGray}}
}

\pagestyle{fancy}
\fancyhf{}
\fancyhead[L]{\small 机器视觉技术练习题合集}
\fancyhead[R]{\small\nouppercase{\leftmark}}
\fancyfoot[C]{\thepage}
\renewcommand{\headrulewidth}{0.45pt}
\renewcommand{\footrulewidth}{0pt}

\newcommand{\EditionName}{\ifshowsolutions 含答案版\else 练习版\fi}
\newcommand{\ChapterIntro}[1]{%
  \begin{tcolorbox}[colback=MVLight,colframe=MVBlue!50,boxrule=0.5pt,arc=1.5mm,left=2mm,right=2mm,top=1.5mm,bottom=1.5mm]
  \small #1
  \end{tcolorbox}
}
\newcommand{\Question}[2]{%
  \par\Needspace{4\baselineskip}\medskip
  \noindent\textbf{\color{MVBlue}#1}\quad #2\par
}
\newenvironment{choices}{\begin{enumerate}[label=\Alph*.,itemsep=0.15em,topsep=0.25em,leftmargin=3em]}{\end{enumerate}}
\newenvironment{AnswerBox}{\begin{tcolorbox}[breakable,colback=MVLight,colframe=MVBlue!45,title={参考答案},fonttitle=\bfseries,boxrule=0.55pt,arc=1.5mm,left=2mm,right=2mm,top=1.5mm,bottom=1.5mm]}{\end{tcolorbox}}
\newcommand{\PracticeSpace}[1]{\ifshowsolutions\else\par\vspace{#1}\par\fi}
\DefineVerbatimEnvironment{CodeBlock}{Verbatim}{fontsize=\small,breaklines=true,breakanywhere=true,frame=single,framesep=3mm,rulecolor=\color{MVBlue!35},formatcom=\color{black}}
\pretocmd{\section}{\Needspace{8\baselineskip}}{}{}

\begin{document}
\frontmatter
\begin{titlepage}
  \centering
  \vspace*{2.0cm}
  {\color{MVBlue}\rule{0.82\textwidth}{1.2pt}}\\[1.3cm]
  {\fontsize{30pt}{38pt}\selectfont\bfseries 机器视觉技术\\[0.35em]练习题合集\par}
  \vspace{0.9cm}
  {\Large\color{MVGold}\EditionName\par}
  \vspace{1.2cm}
  {\large 模拟练习卷 · 深度学习专项 · Faster R-CNN专项\par}
  \vfill
  {\color{MVBlue}\rule{0.82\textwidth}{1.2pt}}\\[0.8cm]
  {\large NKU AI Study\\[0.4em]\today\par}
\end{titlepage}
\tableofcontents
\mainmatter
\chapter{机器视觉技术模拟练习卷 (一)}
\ChapterIntro{本卷满分100分。建议先完成练习，再对照答案复盘知识点与计算过程。}
\section{一、选择题 (10题，每题2分)}
\Question{1.}{一幅大小为 M×N 的灰度图像共有 L 个灰度级，其直方图 h(k) 表示}
\begin{choices}
\item 灰度值为 k 的像素位置
\item 灰度值为 k 的像素数量
\item 图像中所有像素的平均灰度
\item 相邻像素之间的灰度差
\end{choices}
\Question{2.}{以下关于图像采样和量化的说法，正确的是}
\begin{choices}
\item 采样决定灰度级数量，量化决定空间分辨率
\item 采样和量化都只影响空间分辨率
\item 采样决定空间分辨率，量化决定灰度分辨率
\item 量化级数越少，图像灰度变化越平滑
\end{choices}
\Question{3.}{对椒盐噪声具有较好抑制作用，并且能够较好保留边缘的滤波器是}
\begin{choices}
\item 均值滤波器
\item 高斯滤波器
\item 中值滤波器
\item 拉普拉斯滤波器
\end{choices}
\Question{4.}{Canny边缘检测中，非极大值抑制的主要作用是}
\begin{choices}
\item 增强图像对比度
\item 将较宽的边缘细化为单像素边缘
\item 消除所有弱边缘
\item 对图像进行高斯平滑
\end{choices}
\Question{5.}{对二值图像先进行腐蚀，再进行膨胀，得到的形态学运算是}
\begin{choices}
\item 闭运算
\item 开运算
\item 顶帽运算
\item 边界提取
\end{choices}
\Question{6.}{SIFT特征具有尺度不变性的主要原因是}
\begin{choices}
\item 使用了灰度直方图
\item 在高斯尺度空间中检测极值点
\item 使用了非极大值抑制
\item 对描述子进行了二值化
\end{choices}
\Question{7.}{HOG特征在方向直方图中累加的通常是}
\begin{choices}
\item 像素数量
\item 像素灰度值
\item 梯度幅值
\item 像素坐标
\end{choices}
\Question{8.}{在Hough直线检测中，图像空间中的一个边缘点，在参数空间中对应}
\begin{choices}
\item 一个点
\item 一条曲线
\item 一个矩形区域
\item 一条固定直线
\end{choices}
\Question{9.}{输入特征图大小为 32×32×3，使用16个 3×3 卷积核，步长为1，填充为1，则输出特征图大小为}
\begin{choices}
\item 30×30×16
\item 32×32×16
\item 32×32×3
\item 34×34×16
\end{choices}
\Question{10.}{Faster R-CNN中，ROI Pooling的主要作用是}
\begin{choices}
\item 生成候选框
\item 将不同大小的候选区域转换为固定大小的特征
\item 计算Anchor与真实框的IoU
\item 删除重叠候选框
\end{choices}
\section{二、判断题 (5题，每题2分)}
\Question{1.}{均值滤波属于非线性滤波，中值滤波属于线性滤波。}
\par\noindent\hfill（\hspace{2em}）\par
\Question{2.}{Canny边缘检测中的双阈值操作，可以根据强边缘连接情况保留部分弱边缘。}
\par\noindent\hfill（\hspace{2em}）\par
\Question{3.}{图像经过旋转后，SIFT关键点的主方向也会相应改变，因此其描述子仍能保持较好的旋转不变性。}
\par\noindent\hfill（\hspace{2em}）\par
\Question{4.}{HOG特征通常直接对整幅图像的梯度直方图进行一次归一化，不需要划分Cell和Block。}
\par\noindent\hfill（\hspace{2em}）\par
\Question{5.}{ResNet中的残差连接可以缓解深层网络的退化问题，并有利于梯度传播。}
\par\noindent\hfill（\hspace{2em}）\par
\section{三、简答题 (5题，每题6分)}
\Question{1.}{简述线性滤波和非线性滤波的区别，并分别举出两个典型滤波器。}
\PracticeSpace{2.4cm}
\Question{2.}{简述Canny边缘检测的主要步骤。}
\par\noindent\textbf{要求写出以下内容：}
\par\noindent\textbullet\quad 图像平滑；
\par\noindent\textbullet\quad 梯度计算；
\par\noindent\textbullet\quad 非极大值抑制；
\par\noindent\textbullet\quad 双阈值检测；
\par 边缘连接。
\PracticeSpace{2.4cm}
\Question{3.}{简述SIFT特征检测具有尺度不变性和旋转不变性的原因。}
\PracticeSpace{2.4cm}
\Question{4.}{简述HOG特征的生成步骤，并说明为什么方向直方图中累加的是梯度幅值，而不是像素数量。}
\PracticeSpace{2.4cm}
\Question{5.}{简述ResNet残差结构的基本思想，并说明恒等映射为什么有利于深层网络训练。}
\PracticeSpace{2.4cm}
\section{四、解答题 (2题，每题10分)}
\Question{1.}{卷积运算与边缘检测}
\par 给定图像区域
\[
I = \begin{bmatrix}
1 & 1 & 1 \\
1 & 5 & 5 \\
1 & 5 & 9
\end{bmatrix}
\]
\par 以及水平方向Sobel算子
\[
G_x = \begin{bmatrix}
-1 & -2 & -1 \\
0 & 0 & 0 \\
1 & 2 & 1
\end{bmatrix}
\]
\par 垂直方向Sobel算子
\[
G_y = \begin{bmatrix}
-1 & 0 & 1 \\
-2 & 0 & 2 \\
-1 & 0 & 1
\end{bmatrix}
\]
\par\noindent\textbf{回答：}
\par\noindent\textbullet\quad 计算该区域中心位置的 g\_x 和 g\_y；
\par 计算梯度幅值
\begin{center}\small\ttfamily M = sqrt(g\_x\textasciicircum{}2 + g\_y\textasciicircum{}2)\end{center}
\par 计算梯度方向
\begin{center}\small\ttfamily θ = atan2(g\_y, g\_x)\end{center}
\par 说明该位置更可能属于水平方向边缘还是垂直方向边缘。
\PracticeSpace{4.2cm}
\Question{2.}{卷积神经网络尺寸与参数计算}
\par\noindent\textbf{某卷积神经网络的一部分如下：}
\par 输入：32 × 32 × 3
\par\noindent\textbf{Conv1：}
\par 卷积核大小 3 × 3
\par 输出通道数 16
\begin{center}\small\ttfamily stride = 1 padding = 1\end{center}
\par\noindent\textbf{MaxPool：}
\par 窗口大小 2 × 2
\par stride = 2
\par\noindent\textbf{Conv2：}
\par 卷积核大小 3 × 3
\par 输出通道数 32
\begin{center}\small\ttfamily stride = 1 padding = 0\end{center}
\par\noindent\textbf{回答：}
\par\noindent\textbullet\quad Conv1输出特征图尺寸；
\par\noindent\textbullet\quad Conv1可学习参数数量，包含偏置；
\par\noindent\textbullet\quad MaxPool输出特征图尺寸；
\par\noindent\textbullet\quad Conv2输出特征图尺寸；
\par Conv2可学习参数数量，包含偏置。
\par\noindent\textbf{卷积输出尺寸公式为：}
\begin{center}\small\ttfamily H\_out = floor((H\_in + 2P - K)/S) + 1\end{center}
\PracticeSpace{4.2cm}
\section{五、程序题 (2题，每题10分)}
\Question{1.}{OpenCV非线性滤波程序题}
\par 使用C++和OpenCV编写函数，对输入灰度图像进行 3×3 中值滤波。
\par\noindent\textbf{要求：}
\par\noindent\textbullet\quad 不允许直接调用 cv::medianBlur；
\par\noindent\textbullet\quad 对每个非边界像素，取其 3×3 邻域中的9个灰度值；
\par\noindent\textbullet\quad 将9个灰度值排序；
\par\noindent\textbullet\quad 使用中间值作为输出；
\par 边界像素直接复制原图像值。
\par\noindent\textbf{函数形式：}
\begin{CodeBlock}
cv::Mat myMedianFilter(const cv::Mat& src);
\end{CodeBlock}
\PracticeSpace{5.5cm}
\Question{2.}{PyTorch残差模块程序题}
\par 使用PyTorch编写一个ResNet基本残差模块 BasicBlock。
\par\noindent\textbf{要求：}
\par\noindent\textbullet\quad 包含两个 3×3 卷积层；
\par\noindent\textbullet\quad 每个卷积层后使用BatchNorm；
\par\noindent\textbullet\quad 第一个卷积层后使用ReLU；
\par\noindent\textbullet\quad 将卷积结果与输入恒等映射相加；
\par\noindent\textbullet\quad 相加之后再次使用ReLU；
\par 当输入通道和输出通道不同时，使用 1×1 卷积调整维度。
\par\noindent\textbf{函数框架：}
\begin{CodeBlock}
class BasicBlock(nn.Module):
    def __init__(self, in_channels, out_channels, stride=1):
        pass

    def forward(self, x):
        pass
\end{CodeBlock}
\PracticeSpace{5.5cm}
\ifshowsolutions
\clearpage
\section{参考答案}
\section{一、选择题答案}
\par 题号\quad 答案\quad 核心知识点
\par 1\quad B\quad 直方图统计灰度频数
\par 2\quad C\quad 采样决定空间，量化决定灰度
\par 3\quad C\quad 中值滤波抑制椒盐噪声
\par 4\quad B\quad 非极大值抑制细化边缘
\par 5\quad B\quad 先腐蚀后膨胀为开运算
\par 6\quad B\quad 高斯尺度空间
\par 7\quad C\quad 梯度幅值加权投票
\par 8\quad B\quad 一个图像点对应参数空间中的曲线
\par 9\quad B\quad 填充后尺寸不变
\par 10\quad B\quad 不同ROI转为固定尺寸
\section{二、判断题答案}
\par 错误。均值滤波是线性滤波，中值滤波是非线性滤波。
\par 正确。与强边缘连通的弱边缘可以被保留。
\par 正确。SIFT以关键点主方向作为描述子的参考方向。
\par 错误。HOG要划分Cell和Block，并进行局部归一化。
\par 正确。
\section{三、简答题参考要点}
\Question{1.}{线性滤波与非线性滤波}
\par\noindent\textbf{线性滤波的输出是邻域像素的线性加权和，可以用卷积表示，例如：}
\par\noindent\textbullet\quad 均值滤波；
\par 高斯滤波。
\par\noindent\textbf{非线性滤波的输出不能写成邻域像素的固定线性加权和，例如：}
\par\noindent\textbullet\quad 中值滤波；
\par\noindent\textbullet\quad 最大值滤波；
\par 最小值滤波。
\par 中值滤波对椒盐噪声效果好，并且比均值滤波更容易保留边缘。
\Question{2.}{Canny边缘检测步骤}
\par\noindent\textbullet\quad 使用高斯滤波消除噪声；
\par\noindent\textbullet\quad 计算图像在水平和垂直方向上的梯度；
\par\noindent\textbullet\quad 计算梯度幅值和梯度方向；
\par\noindent\textbullet\quad 沿梯度方向进行非极大值抑制，使边缘变细；
\par\noindent\textbullet\quad 使用高、低双阈值区分强边缘、弱边缘和非边缘；
\par 根据连通性保留与强边缘相连的弱边缘。
\Question{3.}{SIFT的不变性}
\par **尺度不变性：**SIFT建立高斯尺度空间，并在不同尺度的DoG图像中检测极值点。同一目标缩放后，仍可在相应尺度找到相同关键点。
\par **旋转不变性：**SIFT根据关键点邻域梯度方向直方图确定主方向，并以主方向作为描述子参考方向。因此图像旋转时，参考方向也随之旋转。
\Question{4.}{HOG生成步骤}
\par\noindent\textbullet\quad 对输入图像进行归一化或Gamma校正；
\par\noindent\textbullet\quad 计算每个像素的梯度幅值和梯度方向；
\par\noindent\textbullet\quad 将图像划分为多个Cell；
\par\noindent\textbullet\quad 在每个Cell中建立梯度方向直方图；
\par\noindent\textbullet\quad 将若干相邻Cell组成Block；
\par\noindent\textbullet\quad 对Block中的特征进行归一化；
\par 连接所有Block特征形成HOG描述子。
\par 方向直方图中使用梯度幅值作为投票权重，是因为梯度越大的像素通常对应越明显的边缘，应当对特征贡献更大。
\Question{5.}{ResNet基本思想}
\par 普通网络直接学习目标映射 H(x)，ResNet改为学习残差：
\par F(x)=H(x)−x
\par\noindent\textbf{于是网络输出为：}
\par y=F(x)+x
\par 当增加的网络层没有学习到有效特征时，只需要使 F(x)≈0，整个模块就接近恒等映射，不会明显破坏原有特征。
\par 跳跃连接还为梯度提供了更直接的传播路径，可以缓解梯度消失和网络退化问题。
\section{四、解答题答案}
\Question{1.}{Sobel计算}
\par 计算 g
\begin{center}\small\ttfamily g\_x = (-1)×1 + 0×1 + 1×1 + (-2)×1 + 0×5 + 2×5 + (-1)×1 + 0×5 + 1×9 = 16\end{center}
\par 计算 g
\begin{center}\small\ttfamily g\_y = (-1)×1 + (-2)×1 + (-1)×1 + 0 + 1×1 + 2×5 + 1×9 = 16\end{center}
\par 梯度幅值
\begin{center}\small\ttfamily M = sqrt(16\textasciicircum{}2 + 16\textasciicircum{}2) = 16 sqrt(2) ≈ 22.63\end{center}
\par 梯度方向
\begin{center}\small\ttfamily θ = atan2(g\_y, g\_x)\end{center}
\par 梯度方向为 45
\par ∘
\par ，说明该位置亮度变化同时包含明显的水平变化和垂直变化，属于斜向边缘，而不是纯水平或纯垂直边缘。
\par 注意：梯度方向与边缘延伸方向垂直。
\Question{2.}{CNN计算}
\par Conv1输出尺寸
\begin{center}\small\ttfamily H\_out = (32 + 2×1 - 3)/1 + 1 = 32\end{center}
\par\noindent\textbf{因此：}
\par 32×32×16
\par Conv1参数量
\par\noindent\textbf{每个卷积核参数：}
\par 3×3×3
\par\noindent\textbf{共有16个卷积核，并且每个卷积核有一个偏置：}
\par (3×3×3+1)×16=28×16=448
\par 池化后尺寸
\par 16×16×16
\par Conv2输出尺寸
\begin{center}\small\ttfamily H\_out = (16 - 3)/1 + 1 = 14\end{center}
\par\noindent\textbf{因此：}
\par 14×14×32
\par Conv2参数量
\begin{center}\small\ttfamily (3×3×16+1)×32 =145×32=4640\end{center}
\section{五、程序题参考答案}
\Question{1.}{手写中值滤波}
\begin{CodeBlock}
#include <opencv2/opencv.hpp>
#include <algorithm>
#include <stdexcept>
#include <vector>

cv::Mat myMedianFilter(const cv::Mat& src)
{
    if (src.empty()) {
        throw std::invalid_argument("输入图像不能为空");
    }

    if (src.type() != CV_8UC1) {
        throw std::invalid_argument("输入图像必须是8位单通道灰度图像");
    }

    cv::Mat dst = src.clone();

    for (int i = 1; i < src.rows - 1; ++i) {
        for (int j = 1; j < src.cols - 1; ++j) {
            std::vector<uchar> values;
            values.reserve(9);

            for (int di = -1; di <= 1; ++di) {
                for (int dj = -1; dj <= 1; ++dj) {
                    values.push_back(
                        src.at<uchar>(i + di, j + dj)
                    );
                }
            }

            std::sort(values.begin(), values.end());

            // 排序后第5个元素，即下标4
            dst.at<uchar>(i, j) = values[4];
        }
    }

    return dst;
}
\end{CodeBlock}
\par int main()
\par \{
\begin{CodeBlock}
    cv::Mat image =
        cv::imread("input.png", cv::IMREAD_GRAYSCALE);

    if (image.empty()) {
        return -1;
    }

    cv::Mat result = myMedianFilter(image);

    cv::imshow("Original", image);
    cv::imshow("Median Filter", result);
    cv::waitKey(0);

    return 0;
}
\end{CodeBlock}
\par\noindent\textbf{核心语句是：}
\begin{CodeBlock}
std::sort(values.begin(), values.end());
dst.at<uchar>(i, j) = values[4];
\end{CodeBlock}
\par 因为9个数字排序后，中值是下标为4的元素。
\Question{2.}{ResNet基本残差模块}
\begin{CodeBlock}
import torch
import torch.nn as nn

class BasicBlock(nn.Module):
    def __init__(
        self,
        in_channels: int,
        out_channels: int,
        stride: int = 1
    ):
        super().__init__()

        self.conv1 = nn.Conv2d(
            in_channels,
            out_channels,
            kernel_size=3,
            stride=stride,
            padding=1,
            bias=False
        )

        self.bn1 = nn.BatchNorm2d(out_channels)
        self.relu = nn.ReLU(inplace=True)

        self.conv2 = nn.Conv2d(
            out_channels,
            out_channels,
            kernel_size=3,
            stride=1,
            padding=1,
            bias=False
        )

        self.bn2 = nn.BatchNorm2d(out_channels)

        # 输入、输出尺寸不一致时调整恒等分支
        if stride != 1 or in_channels != out_channels:
            self.shortcut = nn.Sequential(
                nn.Conv2d(
                    in_channels,
                    out_channels,
                    kernel_size=1,
                    stride=stride,
                    bias=False
                ),
                nn.BatchNorm2d(out_channels)
            )
        else:
            self.shortcut = nn.Identity()

    def forward(self, x: torch.Tensor) -> torch.Tensor:
        identity = self.shortcut(x)

        out = self.conv1(x)
        out = self.bn1(out)
        out = self.relu(out)

        out = self.conv2(out)
        out = self.bn2(out)

        out = out + identity
        out = self.relu(out)

        return out

if __name__ == "__main__":
    x = torch.randn(4, 64, 32, 32)

    block = BasicBlock(
        in_channels=64,
        out_channels=128,
        stride=2
    )

    y = block(x)

    print("输入尺寸：", x.shape)
    print("输出尺寸：", y.shape)
\end{CodeBlock}
\par\noindent\textbf{输出尺寸为：}
\par 输入尺寸：torch.Size([4, 64, 32, 32])
\par 输出尺寸：torch.Size([4, 128, 16, 16])
\par\noindent\textbf{其中 1×1 卷积同时完成：}
\par\noindent\textbullet\quad 通道数由64变为128；
\par\noindent\textbullet\quad 空间尺寸由 32×32 变为 16×16；
\par 保证恒等分支能够与主分支相加。
\fi
\chapter{机器视觉技术模拟练习卷 (二)}
\ChapterIntro{本卷满分100分。建议先完成练习，再对照答案复盘知识点与计算过程。}
\section{一、选择题 (10题，每题2分)}
\Question{1.}{一幅8位灰度图像最多可以表示的灰度级数量为}
\begin{choices}
\item 8
\item 16
\item 255
\item 256
\end{choices}
\Question{2.}{关于图像直方图，下列说法错误的是}
\begin{choices}
\item 直方图能够反映图像中各灰度级的像素数量
\item 直方图可以反映图像整体的明暗分布
\item 直方图完整保留了像素的空间位置信息
\item 两幅不同的图像可能具有相同的直方图
\end{choices}
\Question{3.}{与均值滤波相比，高斯滤波的主要特点是}
\begin{choices}
\item 所有邻域像素权重完全相同
\item 距离中心越近的像素通常权重越大
\item 属于非线性滤波
\item 主要用于检测直线
\end{choices}
\Question{4.}{拉普拉斯算子进行边缘检测时，对噪声比较敏感，主要原因是}
\begin{choices}
\item 拉普拉斯算子只计算图像平均值
\item 拉普拉斯算子属于二阶微分算子
\item 拉普拉斯算子不能用于灰度图像
\item 拉普拉斯算子不包含卷积运算
\end{choices}
\Question{5.}{对二值图像先进行膨胀，再进行腐蚀，称为}
\begin{choices}
\item 开运算
\item 闭运算
\item 边界提取
\item 骨架提取
\end{choices}
\Question{6.}{SIFT算法在检测尺度空间极值点后，需要删除低对比度关键点，主要是为了}
\begin{choices}
\item 减少不稳定关键点
\item 增大图像尺寸
\item 保证所有关键点位于边缘
\item 将彩色图像转换为灰度图像
\end{choices}
\Question{7.}{HOG特征中，对相邻若干Cell组成的Block进行归一化，主要目的是}
\begin{choices}
\item 增加图像的空间尺寸
\item 减小光照和局部对比度变化的影响
\item 将梯度方向变成灰度值
\item 删除所有弱梯度
\end{choices}
\Question{8.}{Hough直线变换通常采用的直线参数方程是}
\begin{choices}
\item y=kx+b
\item x²+y²=r²
\item ρ=xcosθ+ysinθ
\item y=ax²+bx+c
\end{choices}
\Question{9.}{输入特征图大小为 28×28×3，使用6个 5×5 卷积核，步长为1，不填充。输出特征图尺寸和可学习参数量分别为}
\begin{choices}
\item 24×24×6，456
\item 28×28×6，450
\item 24×24×3，456
\item 32×32×6，450
\end{choices}
\par 假设每个卷积核包含一个偏置参数。
\Question{10.}{U-Net中编码器与解码器之间使用跳跃连接的主要作用是}
\begin{choices}
\item 删除浅层特征
\item 融合浅层定位信息和深层语义信息
\item 生成目标检测Anchor
\item 将所有特征图压缩为一维向量
\end{choices}
\section{二、判断题 (5题，每题2分)}
\Question{1.}{图像的曼哈顿距离 D}
\par 4
\par 一定不小于棋盘距离 D
\par 8
\par 。
\par\noindent\hfill（\hspace{2em}）\par
\Question{2.}{高斯滤波可以在完全不损失任何边缘信息的情况下去除所有噪声。}
\par\noindent\hfill（\hspace{2em}）\par
\Question{3.}{对同一二值图像进行膨胀时，结构元素越大，前景区域通常扩张得越明显。}
\par\noindent\hfill（\hspace{2em}）\par
\Question{4.}{SIFT算法中的DoG图像可以由相邻尺度的高斯平滑图像相减得到。}
\par\noindent\hfill（\hspace{2em}）\par
\Question{5.}{Faster R-CNN中的RPN需要判断Anchor是否包含目标，并对候选框位置进行回归。}
\par\noindent\hfill（\hspace{2em}）\par
\section{三、简答题 (5题，每题6分)}
\Question{1.}{简述图像采样和量化的含义，并分别说明采样不足和量化级数不足可能造成的现象。}
\PracticeSpace{2.4cm}
\Question{2.}{为什么在进行边缘检测之前，通常需要先对图像进行平滑处理？}
\par 要求结合微分运算和图像噪声进行说明。
\PracticeSpace{2.4cm}
\Question{3.}{简述二值图像开运算和闭运算的执行顺序及主要作用。}
\PracticeSpace{2.4cm}
\Question{4.}{简述HOG特征的基本生成过程，并说明Cell和Block分别起什么作用。}
\PracticeSpace{2.4cm}
\Question{5.}{简述Faster R-CNN的基本检测流程，并说明Anchor、RPN和ROI Pooling分别承担什么作用。}
\PracticeSpace{2.4cm}
\section{四、解答题 (2题，每题10分)}
\Question{1.}{二值形态学与Hough变换}
\par\noindent\textbf{（1）} 二值形态学计算
\par 给定二值图像
\[
I = \begin{bmatrix}
0 & 0 & 0 & 0 & 0 \\
0 & 1 & 1 & 0 & 0 \\
0 & 1 & 1 & 1 & 0 \\
0 & 0 & 1 & 0 & 0 \\
0 & 0 & 0 & 0 & 0
\end{bmatrix}
\]
\par 结构元素为
\[
B = \begin{bmatrix}
0 & 1 & 0 \\
1 & 1 & 1 \\
0 & 1 & 0
\end{bmatrix}
\]
\par 结构元素的原点位于中心，图像外部像素视为0。
\par\noindent\textbf{计算：}
\par\noindent\textbullet\quad 图像 I 被 B 膨胀后的结果；
\par 图像 I 被 B 腐蚀后的结果。
\par\noindent\textbf{（2）} Hough直线检测
\par\noindent\textbf{图像中存在以下6个边缘点：}
\par (1,1),(2,2),(3,3),(2,0),(2,1),(2,3)
\par 采用Hough直线方程
\par ρ=xcosθ+ysinθ
\par\noindent\textbf{只考虑以下两个角度：}
\begin{center}\small\ttfamily θ=0° ,θ=135°\end{center}
\par\noindent\textbf{回答：}
\par 当 θ=0
\par ∘
\par 时，哪个 ρ 值获得的票数最多？
\par 当 θ=135
\par ∘
\par 时，哪个 ρ 值获得的票数最多？
\par 写出检测到的两条直线。
\PracticeSpace{4.2cm}
\Question{2.}{CNN特征尺寸与参数计算}
\par\noindent\textbf{某卷积神经网络结构如下：}
\par 输入图像：64 × 64 × 3
\par\noindent\textbf{Conv1：}
\par 卷积核大小：5 × 5
\par 输出通道数：8
\begin{center}\small\ttfamily stride = 2 padding = 2\end{center}
\par\noindent\textbf{MaxPool：}
\par 窗口大小：2 × 2
\par stride = 2
\par\noindent\textbf{Conv2：}
\par 卷积核大小：3 × 3
\par 输出通道数：16
\begin{center}\small\ttfamily stride = 1 padding = 1\end{center}
\par Flatten
\par\noindent\textbf{全连接层：}
\par 输出神经元数量为10
\par\noindent\textbf{回答：}
\par\noindent\textbullet\quad Conv1输出特征图尺寸；
\par\noindent\textbullet\quad Conv1参数量，包含偏置；
\par\noindent\textbullet\quad MaxPool输出特征图尺寸；
\par\noindent\textbullet\quad Conv2输出特征图尺寸；
\par\noindent\textbullet\quad Conv2参数量，包含偏置；
\par\noindent\textbullet\quad 全连接层参数量；
\par 整个网络中上述可学习参数总数。
\par\noindent\textbf{卷积输出尺寸公式为：}
\begin{center}\small\ttfamily H\_out = floor((H\_in + 2P - K)/S) + 1\end{center}
\PracticeSpace{4.2cm}
\section{五、程序题 (2题，每题10分)}
\Question{1.}{OpenCV最大值滤波}
\par 使用C++和OpenCV编写一个 3×3 最大值滤波函数。
\par 最大值滤波的定义为：对于每个像素，取其 3×3 邻域内的最大灰度值作为输出。
\par\noindent\textbf{要求：}
\par\noindent\textbullet\quad 不允许调用 cv::dilate；
\par\noindent\textbullet\quad 不允许调用其他现成最大值滤波函数；
\par\noindent\textbullet\quad 输入为8位单通道灰度图像；
\par\noindent\textbullet\quad 对非边界像素遍历 3×3 邻域；
\par 边界像素直接复制原图像值。
\par\noindent\textbf{函数形式如下：}
\begin{CodeBlock}
cv::Mat maxFilter3x3(const cv::Mat& src);
\end{CodeBlock}
\PracticeSpace{5.5cm}
\Question{2.}{PyTorch简化U-Net}
\par 使用PyTorch编写一个简化的U-Net网络。
\par\noindent\textbf{网络结构如下：}
\par 输入图像
\begin{CodeBlock}
    ↓
\end{CodeBlock}
\par DoubleConv：3 → 64
\begin{CodeBlock}
    ↓
\end{CodeBlock}
\par 最大池化
\begin{CodeBlock}
    ↓
\end{CodeBlock}
\par DoubleConv：64 → 128
\begin{CodeBlock}
    ↓
\end{CodeBlock}
\par 转置卷积上采样：128 → 64
\begin{CodeBlock}
    ↓
\end{CodeBlock}
\par 与编码器对应特征拼接
\begin{CodeBlock}
    ↓
\end{CodeBlock}
\par DoubleConv：128 → 64
\begin{CodeBlock}
    ↓
\end{CodeBlock}
\par 1 × 1卷积输出分类结果
\par\noindent\textbf{要求：}
\par\noindent\textbullet\quad DoubleConv包含两个 3×3 卷积；
\par\noindent\textbullet\quad 每个卷积后使用BatchNorm和ReLU；
\par\noindent\textbullet\quad 使用最大池化完成下采样；
\par\noindent\textbullet\quad 使用转置卷积完成上采样；
\par\noindent\textbullet\quad 使用 torch.cat 实现跳跃连接；
\par 最后一层使用 1×1 卷积输出 num\_classes 个通道。
\PracticeSpace{5.5cm}
\ifshowsolutions
\clearpage
\section{参考答案}
\section{一、选择题答案}
\par 题号\quad 答案\quad 知识点
\par 1\quad D\quad 8位图像有 2
\par 8
\par =256 个灰度级
\par 2\quad C\quad 直方图不包含空间位置信息
\par 3\quad B\quad 高斯核中心权重大
\par 4\quad B\quad 二阶微分对噪声敏感
\par 5\quad B\quad 先膨胀后腐蚀是闭运算
\par 6\quad A\quad 删除低对比度不稳定点
\par 7\quad B\quad 减小光照和对比度影响
\par 8\quad C\quad Hough直线极坐标形式
\par 9\quad A\quad 尺寸与参数量计算
\par 10\quad B\quad 融合浅层与深层特征
\par 第9题计算
\par\noindent\textbf{输出尺寸：}
\begin{center}\small\ttfamily H\_out = floor((H\_in + 2P - K)/S) + 1\end{center}
\par\noindent\textbf{所以输出为：}
\par 24×24×6
\par\noindent\textbf{参数量：}
\begin{center}\small\ttfamily (5×5×3+1)×6 =76×6=456\end{center}
\section{二、判断题答案}
\par 正确
\par\noindent\textbf{对于两个像素间坐标差 (Δx,Δy)：}
\begin{center}\small\ttfamily D\_4\end{center}
\par =∣Δx∣+∣Δy∣
\begin{center}\small\ttfamily D\_8\end{center}
\par =max(∣Δx∣,∣Δy∣)
\par\noindent\textbf{因此：}
\begin{center}\small\ttfamily D\_4 ≥D\_8\end{center}
\par 错误。平滑噪声的同时通常也会使部分边缘变模糊。
\par 正确。较大的结构元素通常会使前景扩张范围更大。
\par 正确。DoG是相邻高斯尺度图像之差。
\par 正确。RPN包含前景与背景分类以及边界框回归。
\section{三、简答题参考答案}
\Question{1.}{采样和量化}
\par 采样是将连续图像的空间坐标离散化，决定图像的空间分辨率。
\par\noindent\textbf{采样点过少可能出现：}
\par\noindent\textbullet\quad 图像细节丢失；
\par\noindent\textbullet\quad 锯齿；
\par 混叠现象。
\par 量化是将连续的亮度或颜色值离散为有限个等级，决定灰度分辨率。
\par\noindent\textbf{量化级数过少可能出现：}
\par\noindent\textbullet\quad 灰度层次减少；
\par\noindent\textbullet\quad 平滑区域出现假轮廓；
\par 图像色阶不连续。
\Question{2.}{边缘检测前进行平滑的原因}
\par 边缘检测通常需要进行一阶或二阶微分，而微分运算会放大图像中的高频成分。
\par 图像噪声通常也属于高频成分，因此直接进行微分可能产生大量虚假边缘。
\par\noindent\textbf{先采用高斯滤波等方法进行平滑，可以：}
\par\noindent\textbullet\quad 抑制噪声；
\par\noindent\textbullet\quad 减少虚假梯度；
\par 提高边缘检测的稳定性。
\par 但是平滑过强也会使真实边缘变得模糊，因此需要在去噪和保留边缘之间进行权衡。
\Question{3.}{开运算和闭运算}
\par 开运算
\par\noindent\textbf{先腐蚀，后膨胀：}
\par A∘B=(A⊖B)⊕B
\par\noindent\textbf{主要作用：}
\par\noindent\textbullet\quad 删除较小的前景噪声；
\par\noindent\textbullet\quad 断开细小连接；
\par 平滑物体外边界。
\par 闭运算
\par\noindent\textbf{先膨胀，后腐蚀：}
\par A∙B=(A⊕B)⊖B
\par\noindent\textbf{主要作用：}
\par\noindent\textbullet\quad 填补较小的内部孔洞；
\par\noindent\textbullet\quad 连接相邻目标；
\par 弥合狭窄断裂。
\Question{4.}{HOG特征生成过程}
\par\noindent\textbullet\quad 对图像进行灰度化或Gamma归一化；
\par\noindent\textbullet\quad 计算每个像素的水平和垂直梯度；
\par\noindent\textbullet\quad 得到梯度幅值和梯度方向；
\par\noindent\textbullet\quad 将图像划分成多个Cell；
\par\noindent\textbullet\quad 在每个Cell中建立梯度方向直方图；
\par\noindent\textbullet\quad 将相邻Cell组成Block；
\par\noindent\textbullet\quad 对Block中的特征向量进行归一化；
\par 拼接所有Block特征得到最终HOG描述子。
\par Cell的作用
\par Cell用于统计局部区域内的边缘方向分布。
\par Block的作用
\par Block对相邻Cell的特征进行局部归一化，从而减小光照、阴影和局部对比度变化的影响。
\Question{5.}{Faster R-CNN基本流程}
\par\noindent\textbullet\quad 输入图像经过卷积网络，得到共享特征图；
\par\noindent\textbullet\quad RPN在特征图上生成并筛选候选区域；
\par\noindent\textbullet\quad 将候选区域映射到共享特征图；
\par\noindent\textbullet\quad ROI Pooling将不同大小的候选区域转换为固定大小的特征；
\par\noindent\textbullet\quad 全连接网络对候选区域进行分类和边界框回归；
\par 使用非极大值抑制得到最终检测结果。
\par Anchor
\par Anchor是在特征图每个位置预设的不同比例和尺寸的参考框。
\par RPN
\par\noindent\textbf{RPN负责：}
\par\noindent\textbullet\quad 判断Anchor是前景还是背景；
\par\noindent\textbullet\quad 预测Anchor到真实目标框的坐标偏移；
\par 生成候选区域。
\par ROI Pooling
\par ROI Pooling将不同大小的候选区域变成固定大小的特征，以便输入后续全连接层。
\section{四、解答题参考答案}
\Question{1.}{二值形态学与Hough变换}
\par\noindent\textbf{（1）} 膨胀结果
\par 结构元素包含中心及其上、下、左、右位置。
\par\noindent\textbf{每个前景像素都会向四个方向扩张一个像素，因此：}
\[
I\oplus B = \begin{bmatrix}
0 & 1 & 1 & 0 & 0 \\
1 & 1 & 1 & 1 & 0 \\
1 & 1 & 1 & 1 & 1 \\
0 & 1 & 1 & 1 & 0 \\
0 & 0 & 1 & 0 & 0
\end{bmatrix}
\]
\par 腐蚀结果
\par 只有当某个像素及其上、下、左、右位置全部为1时，该像素才会在腐蚀结果中保留。
\par\noindent\textbf{因此：}
\[
I\ominus B = \begin{bmatrix}
0 & 0 & 0 & 0 & 0 \\
0 & 0 & 0 & 0 & 0 \\
0 & 0 & 1 & 0 & 0 \\
0 & 0 & 0 & 0 & 0 \\
0 & 0 & 0 & 0 & 0
\end{bmatrix}
\]
\par\noindent\textbf{（2）} Hough变换
\par 当 θ=0
\par ∘
\par\noindent\textbf{因为：}
\begin{center}\small\ttfamily cos0° =1,sin0° =0\end{center}
\par\noindent\textbf{所以：}
\par ρ=x
\par\noindent\textbf{各点对应的 ρ 为：}
\par 1,2,3,2,2,2
\par\noindent\textbf{其中：}
\par ρ=2
\par 获得4票，票数最多。
\par\noindent\textbf{对应直线：}
\par x=2
\par 当 θ=135
\begin{center}\small\ttfamily cos135° = -sqrt(2)/2，sin135° = sqrt(2)/2\end{center}
\par\noindent\textbf{因此：}
\begin{center}\small\ttfamily ρ = (sqrt(2)/2)(y - x)\end{center}
\par\noindent\textbf{对于点：}
\par (1,1),(2,2),(3,3)
\par\noindent\textbf{均有：}
\par ρ=0
\par 因此 ρ=0 获得3票。
\par\noindent\textbf{对应直线：}
\par 0=−x+y
\par\noindent\textbf{即：}
\par y=x
\par\noindent\textbf{最终检测到的两条直线为：}
\begin{center}\small\ttfamily x=2 y=x\end{center}
\Question{2.}{CNN尺寸和参数计算}
\par\noindent\textbf{（1）} Conv1输出尺寸
\begin{center}\small\ttfamily H\_out = floor((64 + 2×2 - 5)/2) + 1 = 32\end{center}
\par\noindent\textbf{因此输出特征图为：}
\par 32×32×8
\par\noindent\textbf{（2）} Conv1参数量
\par\noindent\textbf{每个卷积核连接3个输入通道：}
\par 5×5×3=75
\par 加一个偏置，共76个参数。
\par\noindent\textbf{一共有8个卷积核：}
\begin{center}\small\ttfamily (5×5×3+1)×8 = 76×8 = 608\end{center}
\par\noindent\textbf{（3）} MaxPool输出尺寸
\par\noindent\textbf{池化窗口为 2×2，步长为2：}
\par 16×16×8
\par\noindent\textbf{（4）} Conv2输出尺寸
\par\noindent\textbf{输入尺寸为 16×16，卷积核为3，填充为1，步长为1：}
\begin{center}\small\ttfamily H\_out = (16 + 2×1 - 3)/1 + 1 = 16\end{center}
\par\noindent\textbf{因此：}
\par 16×16×16
\par\noindent\textbf{（5）} Conv2参数量
\par\noindent\textbf{每个卷积核参数量为：}
\par 3×3×8+1=73
\par\noindent\textbf{共有16个卷积核：}
\begin{center}\small\ttfamily 73×16 = 1168\end{center}
\par\noindent\textbf{（6）} 全连接层参数量
\par\noindent\textbf{Flatten后的特征数量为：}
\par 16×16×16=4096
\par\noindent\textbf{全连接层有10个输出：}
\begin{center}\small\ttfamily 4096×10 + 10 = 40970\end{center}
\par\noindent\textbf{（7）} 总参数量
\par 608+1168+40970=42746
\par\noindent\textbf{因此总参数量为：}
\par 42746
\section{五、程序题参考答案}
\Question{1.}{手写 3×3 最大值滤波}
\begin{CodeBlock}
#include <opencv2/opencv.hpp>
#include <algorithm>
#include <stdexcept>

cv::Mat maxFilter3x3(const cv::Mat& src)
{
    if (src.empty()) {
        throw std::invalid_argument("输入图像不能为空");
    }

    if (src.type() != CV_8UC1) {
        throw std::invalid_argument(
            "输入图像必须是8位单通道灰度图像"
        );
    }

    // 边界默认复制原图像
    cv::Mat dst = src.clone();

    for (int row = 1; row < src.rows - 1; ++row) {
        for (int col = 1; col < src.cols - 1; ++col) {

            uchar maxValue = 0;

            // 遍历3×3邻域
            for (int dr = -1; dr <= 1; ++dr) {
                for (int dc = -1; dc <= 1; ++dc) {
                    uchar value =
                        src.at<uchar>(row + dr, col + dc);

                    maxValue = std::max(maxValue, value);
                }
            }

            dst.at<uchar>(row, col) = maxValue;
        }
    }

    return dst;
}
\end{CodeBlock}
\par int main()
\par \{
\begin{CodeBlock}
    cv::Mat src = cv::imread(
        "input.png",
        cv::IMREAD_GRAYSCALE
    );

    if (src.empty()) {
        return -1;
    }

    cv::Mat dst = maxFilter3x3(src);

    cv::imshow("Original", src);
    cv::imshow("Maximum Filter", dst);
    cv::waitKey(0);

    return 0;
}
\end{CodeBlock}
\par 核心逻辑
\begin{CodeBlock}
uchar maxValue = 0;

for (int dr = -1; dr <= 1; ++dr) {
    for (int dc = -1; dc <= 1; ++dc) {
        maxValue = std::max(
            maxValue,
            src.at<uchar>(row + dr, col + dc)
        );
    }
}
\end{CodeBlock}
\par 最大值滤波相当于灰度图像的膨胀运算，能够扩大图像中的亮区域。
\Question{2.}{简化U-Net实现}
\begin{CodeBlock}
import torch
import torch.nn as nn

class DoubleConv(nn.Module):
    """连续执行两次卷积、BN和ReLU。"""

    def __init__(self, in_channels: int, out_channels: int):
        super().__init__()

        self.layers = nn.Sequential(
            nn.Conv2d(
                in_channels,
                out_channels,
                kernel_size=3,
                padding=1,
                bias=False
            ),
            nn.BatchNorm2d(out_channels),
            nn.ReLU(inplace=True),

            nn.Conv2d(
                out_channels,
                out_channels,
                kernel_size=3,
                padding=1,
                bias=False
            ),
            nn.BatchNorm2d(out_channels),
            nn.ReLU(inplace=True)
        )

    def forward(self, x: torch.Tensor) -> torch.Tensor:
        return self.layers(x)

class MiniUNet(nn.Module):
    def __init__(
        self,
        in_channels: int = 3,
        num_classes: int = 2
    ):
        super().__init__()

        # 编码器第一层
        self.encoder1 = DoubleConv(
            in_channels,
            64
        )

        self.pool = nn.MaxPool2d(
            kernel_size=2,
            stride=2
        )

        # 编码器第二层
        self.encoder2 = DoubleConv(
            64,
            128
        )

        # 上采样：通道128变为64
        self.up = nn.ConvTranspose2d(
            128,
            64,
            kernel_size=2,
            stride=2
        )

        # 拼接后通道数为64 + 64 = 128
        self.decoder1 = DoubleConv(
            128,
            64
        )

        # 输出每个像素的类别分数
        self.output_conv = nn.Conv2d(
            64,
            num_classes,
            kernel_size=1
        )

    def forward(self, x: torch.Tensor) -> torch.Tensor:
        # 编码器浅层特征
        feature1 = self.encoder1(x)

        # 下采样
        x = self.pool(feature1)

        # 深层特征
        x = self.encoder2(x)

        # 上采样
        x = self.up(x)

        # 跳跃连接：沿通道维拼接
        x = torch.cat(
            [feature1, x],
            dim=1
        )

        # 解码器
        x = self.decoder1(x)

        # 输出分割结果
        x = self.output_conv(x)

        return x

if __name__ == "__main__":
    model = MiniUNet(
        in_channels=3,
        num_classes=4
    )

    image = torch.randn(
        2,
        3,
        256,
        256
    )

    output = model(image)

    print("输入尺寸：", image.shape)
    print("输出尺寸：", output.shape)
\end{CodeBlock}
\par\noindent\textbf{运行结果为：}
\par 输入尺寸： torch.Size([2, 3, 256, 256])
\par 输出尺寸： torch.Size([2, 4, 256, 256])
\par 输出中的4个通道表示每个像素属于4个不同类别的预测分数。U-Net没有直接使用全连接层，因此能够保留图像的空间结构，实现逐像素分类。
\fi
\chapter{深度学习专项题库}
\ChapterIntro{覆盖CNN组成与训练、AlexNet、VGG、GoogLeNet、ResNet、Faster R-CNN、Anchor、ROI Pooling和U-Net。}
\section{一、问答题专项}
\Question{1.}{简述卷积神经网络的基本组成及各部分作用}
\ifshowsolutions
\begin{AnswerBox}
\par\noindent\textbf{卷积神经网络通常由以下部分组成：}
\par 卷积层：通过卷积核提取局部特征，例如边缘、纹理和高级语义特征。
\par 激活函数：引入非线性，使网络能够学习复杂映射，常用ReLU。
\par 池化层：降低特征图空间尺寸，减少计算量，提高对局部位移的鲁棒性。
\par 归一化层：稳定特征分布，加快网络训练，例如Batch Normalization。
\par 全连接层：将提取的特征映射为最终类别或回归结果。
\par 输出层：根据任务输出类别概率、边界框或像素分类结果。
\end{AnswerBox}
\else
\PracticeSpace{2.4cm}
\fi
\Question{2.}{卷积神经网络为什么采用局部连接和权值共享？}
\ifshowsolutions
\begin{AnswerBox}
\par 局部连接是指一个输出神经元只与输入图像的局部区域连接，符合图像局部像素相关性较强的特点。
\par 权值共享是指同一个卷积核在整幅图像上重复使用，使其能够在不同位置检测同一种特征。
\par\noindent\textbf{二者的作用是：}
\par\noindent\textbullet\quad 大幅减少参数量；
\par\noindent\textbullet\quad 降低计算复杂度；
\par\noindent\textbullet\quad 能够检测不同位置的相同特征；
\par 使卷积具有一定的平移等变性。
\end{AnswerBox}
\else
\PracticeSpace{2.4cm}
\fi
\Question{3.}{卷积层与全连接层有什么区别？}
\ifshowsolutions
\begin{AnswerBox}
\par 对比项\quad 卷积层\quad 全连接层
\par 连接方式\quad 局部连接\quad 所有输入与所有输出连接
\par 权值使用\quad 权值共享\quad 每条连接通常有独立权重
\par 参数量\quad 较少\quad 通常较多
\par 是否保留空间结构\quad 保留\quad 展平后通常不保留
\par 主要作用\quad 提取局部特征\quad 完成分类或回归
\par 卷积层适合处理图像等具有空间结构的数据；全连接层通常用于最终决策。
\end{AnswerBox}
\else
\PracticeSpace{2.4cm}
\fi
\Question{4.}{如何计算卷积层输出尺寸和参数量？}
\ifshowsolutions
\begin{AnswerBox}
\par 输入特征图空间尺寸为 H
\begin{center}\small\ttfamily H\_in × W\_in\end{center}
\par\noindent\textbf{，卷积核大小为 K，填充为 P，步长为 S，则：}
\begin{center}\small\ttfamily H\_out = floor((H\_in + 2P - K)/S) + 1\end{center}
\par 宽度计算方法相同。
\par 如果输入通道数为 C
\par in
\par ，输出通道数为 C
\par out
\par\noindent\textbf{，包含偏置时参数量为：}
\begin{center}\small\ttfamily N = (K\_h K\_w C\_in + 1) C\_out\end{center}
\par\noindent\textbf{不包含偏置时为：}
\begin{center}\small\ttfamily N = K\_h K\_w C\_in C\_out\end{center}
\par 卷积层参数量与输入图像的高、宽无关。
\end{AnswerBox}
\else
\PracticeSpace{2.4cm}
\fi
\Question{5.}{池化层有什么作用？最大池化和平均池化有什么区别？}
\ifshowsolutions
\begin{AnswerBox}
\par\noindent\textbf{池化层主要作用是：}
\par\noindent\textbullet\quad 减小特征图空间尺寸；
\par\noindent\textbullet\quad 降低计算量和存储量；
\par\noindent\textbullet\quad 扩大后续神经元的有效感受野；
\par 提高对局部平移和形变的鲁棒性。
\par 最大池化取邻域最大值，更强调显著特征。
\par 平均池化取邻域平均值，更强调区域整体信息。
\par 普通池化层一般没有可学习参数，也通常不改变通道数。
\end{AnswerBox}
\else
\PracticeSpace{2.4cm}
\fi
\Question{6.}{为什么神经网络需要激活函数？}
\ifshowsolutions
\begin{AnswerBox}
\par 如果网络中只有线性变换，那么多层线性变换仍可合并成一次线性变换，增加层数并不能提高网络表达复杂非线性关系的能力。
\par 激活函数为网络引入非线性，使其能够学习复杂的分类边界和图像特征。
\par\noindent\textbf{常见激活函数有：}
\par\noindent\textbullet\quad ReLU；
\par\noindent\textbullet\quad Sigmoid；
\par\noindent\textbullet\quad Tanh；
\par Softmax。
\end{AnswerBox}
\else
\PracticeSpace{2.4cm}
\fi
\Question{7.}{ReLU激活函数有什么特点？}
\ifshowsolutions
\begin{AnswerBox}
\par\noindent\textbf{ReLU定义为：}
\par f(x)=max(0,x)
\par\noindent\textbf{优点：}
\par\noindent\textbullet\quad 计算简单；
\par\noindent\textbullet\quad 正区间梯度为1；
\par\noindent\textbullet\quad 可以缓解梯度消失；
\par\noindent\textbullet\quad 能产生稀疏激活；
\par 通常比Sigmoid收敛更快。
\par\noindent\textbf{缺点：}
\par\noindent\textbullet\quad 当输入长期为负时，输出和梯度均为0；
\par 可能出现“神经元死亡”问题。
\end{AnswerBox}
\else
\PracticeSpace{2.4cm}
\fi
\Question{8.}{简述卷积神经网络的训练循环}
\ifshowsolutions
\begin{AnswerBox}
\par\noindent\textbf{一个标准训练循环包括：}
\par\noindent\textbullet\quad 从数据集中读取一个批次的数据和标签；
\par\noindent\textbullet\quad 将数据输入网络，完成前向传播；
\par\noindent\textbullet\quad 根据预测结果和真实标签计算损失；
\par\noindent\textbullet\quad 清空上一轮保留的梯度；
\par\noindent\textbullet\quad 反向传播计算各参数梯度；
\par\noindent\textbullet\quad 优化器根据梯度更新网络参数；
\par\noindent\textbullet\quad 重复多个批次和多个Epoch；
\par 在验证集上评估模型的泛化性能。
\par\noindent\textbf{PyTorch中的核心顺序为：}
\begin{CodeBlock}
optimizer.zero_grad()
outputs = model(inputs)
loss = criterion(outputs, labels)
\end{CodeBlock}
\begin{center}\small\ttfamily loss.backward() optimizer.step()\end{center}
\end{AnswerBox}
\else
\PracticeSpace{2.4cm}
\fi
\Question{9.}{什么是过拟合？如何缓解？}
\ifshowsolutions
\begin{AnswerBox}
\par 过拟合是指模型在训练集上表现很好，但在验证集或测试集上表现较差，说明模型记住了训练数据中的细节和噪声，泛化能力不足。
\par\noindent\textbf{常用缓解方法：}
\par\noindent\textbullet\quad 增加训练数据；
\par\noindent\textbullet\quad 数据增强；
\par\noindent\textbullet\quad Dropout；
\par\noindent\textbullet\quad 权重衰减；
\par\noindent\textbullet\quad 早停；
\par\noindent\textbullet\quad 降低模型复杂度；
\par\noindent\textbullet\quad Batch Normalization；
\par 使用预训练模型。
\end{AnswerBox}
\else
\PracticeSpace{2.4cm}
\fi
\Question{10.}{简述AlexNet的主要特点}
\ifshowsolutions
\begin{AnswerBox}
\par\noindent\textbf{AlexNet的主要特点包括：}
\par\noindent\textbullet\quad 使用ReLU替代传统Sigmoid或Tanh；
\par\noindent\textbullet\quad 使用Dropout缓解全连接层过拟合；
\par\noindent\textbullet\quad 使用数据增强；
\par\noindent\textbullet\quad 使用最大池化；
\par 采用较深的卷积神经网络完成大规模图像分类。
\par AlexNet证明了深度卷积网络在大规模视觉任务上的有效性。
\end{AnswerBox}
\else
\PracticeSpace{2.4cm}
\fi
\Question{11.}{简述VGGNet的主要特点}
\ifshowsolutions
\begin{AnswerBox}
\par VGGNet的核心特点是大量堆叠 3×3 小卷积核。
\par\noindent\textbf{多个小卷积核堆叠可以：}
\par\noindent\textbullet\quad 获得与大卷积核相近的感受野；
\par\noindent\textbullet\quad 引入更多非线性激活；
\par\noindent\textbullet\quad 减少部分参数；
\par 使网络结构更加规则。
\par 缺点是网络层数较多，全连接层参数量大，计算和存储成本较高。
\end{AnswerBox}
\else
\PracticeSpace{2.4cm}
\fi
\Question{12.}{简述GoogLeNet中的Inception结构}
\ifshowsolutions
\begin{AnswerBox}
\par\noindent\textbf{Inception结构在同一层并行使用不同尺度的操作，例如：}
\par\noindent\textbullet\quad 1×1 卷积；
\par\noindent\textbullet\quad 3×3 卷积；
\par\noindent\textbullet\quad 5×5 卷积；
\par 池化。
\par 然后将各分支输出沿通道维拼接。
\par 其作用是同时提取不同尺度的图像特征。
\par\noindent\textbf{其中 1×1 卷积常用于：}
\par\noindent\textbullet\quad 调整通道数；
\par\noindent\textbullet\quad 降低后续大卷积的计算量；
\par 引入非线性。
\end{AnswerBox}
\else
\PracticeSpace{2.4cm}
\fi
\Question{13.}{简述ResNet残差学习的基本思想}
\ifshowsolutions
\begin{AnswerBox}
\par 普通网络直接学习目标映射 H(x)，ResNet改为学习残差：
\par F(x)=H(x)−x
\par\noindent\textbf{因此输出为：}
\par y=F(x)+x
\par\noindent\textbf{如果新增网络层没有学到有效映射，只需使：}
\par F(x)≈0
\par 残差块就近似恒等映射，不会明显破坏原有特征。
\par\noindent\textbf{残差连接还为梯度提供了直接传播路径，可以缓解：}
\par\noindent\textbullet\quad 梯度消失；
\par\noindent\textbullet\quad 深层网络退化；
\par 深层网络难以优化的问题。
\end{AnswerBox}
\else
\PracticeSpace{2.4cm}
\fi
\Question{14.}{什么是深层网络退化问题？它与过拟合是否相同？}
\ifshowsolutions
\begin{AnswerBox}
\par 网络退化是指随着网络深度增加，训练误差反而升高，深层网络甚至不如浅层网络。
\par 这并不一定是过拟合，因为退化问题中深层网络在训练集上的误差也可能更高。
\par ResNet通过恒等映射和残差学习，使新增层至少能够接近恒等映射，从而缓解网络退化问题。
\end{AnswerBox}
\else
\PracticeSpace{2.4cm}
\fi
\Question{15.}{ResNet中主分支和捷径分支尺寸不一致时怎么办？}
\ifshowsolutions
\begin{AnswerBox}
\par 残差相加要求两条分支的形状完全相同。
\par\noindent\textbf{如果通道数或空间尺寸不同，可以在捷径分支使用：}
\par 1×1
\par\noindent\textbf{卷积，并设置适当的步长，实现：}
\par\noindent\textbullet\quad 通道数调整；
\par\noindent\textbullet\quad 特征图下采样；
\par 尺寸匹配。
\par 之后再进行逐元素相加。
\end{AnswerBox}
\else
\PracticeSpace{2.4cm}
\fi
\Question{16.}{图像分类、目标检测和图像分割有什么区别？}
\ifshowsolutions
\begin{AnswerBox}
\par\noindent\textbf{图像分类：}
\par 判断整幅图像属于哪个类别，只输出类别。
\par\noindent\textbf{目标检测：}
\par 同时输出图像中目标的类别和边界框位置。
\par\noindent\textbf{图像分割：}
\par 对每个像素进行分类，输出与输入图像空间对应的像素级结果。
\par\noindent\textbf{因此三者输出信息从粗到细为：}
\par 整图类别→类别＋位置→逐像素类别
\end{AnswerBox}
\else
\PracticeSpace{2.4cm}
\fi
\Question{17.}{简述R-CNN、Fast R-CNN和Faster R-CNN的区别}
\ifshowsolutions
\begin{AnswerBox}
\par\noindent\textbf{R-CNN：}
\par\noindent\textbullet\quad 使用外部算法生成候选区域；
\par\noindent\textbullet\quad 每个候选区域分别经过CNN；
\par 存在大量重复卷积，速度慢。
\par\noindent\textbf{Fast R-CNN：}
\par\noindent\textbullet\quad 整幅图像只进行一次卷积；
\par\noindent\textbullet\quad 候选区域映射到共享特征图；
\par\noindent\textbullet\quad 使用ROI Pooling获得固定尺寸特征；
\par 仍依赖外部候选框算法。
\par\noindent\textbf{Faster R-CNN：}
\par\noindent\textbullet\quad 使用RPN在共享特征图上生成候选区域；
\par\noindent\textbullet\quad 候选区域生成和检测共享卷积特征；
\par 实现近似端到端训练。
\end{AnswerBox}
\else
\PracticeSpace{2.4cm}
\fi
\Question{18.}{什么是Anchor？它有什么作用？}
\ifshowsolutions
\begin{AnswerBox}
\par Anchor是在特征图每个位置预先设置的一组参考框，通常具有不同的：
\par\noindent\textbullet\quad 尺寸；
\par 长宽比。
\par\noindent\textbf{对于每个Anchor，RPN预测：}
\par\noindent\textbullet\quad 它是前景还是背景；
\par 它相对于真实框的坐标偏移。
\par Anchor为不同大小、不同形状的目标提供边界框回归的初始参考。
\end{AnswerBox}
\else
\PracticeSpace{2.4cm}
\fi
\Question{19.}{简述RPN的作用}
\ifshowsolutions
\begin{AnswerBox}
\par RPN即区域建议网络，它在主干网络输出的共享特征图上滑动，对每个Anchor进行：
\par\noindent\textbullet\quad 前景与背景分类；
\par 边界框位置回归。
\par 之后根据置信度和NMS筛选候选区域，提供给后续检测头。
\par\noindent\textbf{因此RPN的核心输出包括：}
\par\noindent\textbullet\quad 目标性得分；
\par\noindent\textbullet\quad 边界框偏移量；
\par 候选区域Proposal。
\end{AnswerBox}
\else
\PracticeSpace{2.4cm}
\fi
\Question{20.}{ROI Pooling的作用和基本原理是什么？}
\ifshowsolutions
\begin{AnswerBox}
\par 不同候选区域的空间大小不同，而全连接层需要固定长度的输入。
\par ROI Pooling将每个候选区域划分为固定数量的子区域，例如 7×7，并在每个子区域内进行最大池化，最终得到固定大小的特征图。
\par\noindent\textbf{作用是：}
\par\noindent\textbullet\quad 将不同大小的ROI转换为固定尺寸；
\par 使其能够输入后续全连接分类和回归网络。
\par ROI Pooling需要进行坐标取整，因此可能产生空间量化误差。
\end{AnswerBox}
\else
\PracticeSpace{2.4cm}
\fi
\Question{21.}{简述U-Net的结构和工作过程}
\ifshowsolutions
\begin{AnswerBox}
\par U-Net由编码器、解码器和跳跃连接组成。
\par\noindent\textbf{编码器：}
\par\noindent\textbullet\quad 使用卷积提取特征；
\par\noindent\textbullet\quad 使用池化降低空间尺寸；
\par 逐渐获得高级语义信息。
\par\noindent\textbf{解码器：}
\par\noindent\textbullet\quad 使用上采样或转置卷积恢复空间尺寸；
\par 逐步生成像素级分类结果。
\par\noindent\textbf{跳跃连接：}
\par 将编码器对应层的浅层特征与解码器特征拼接，融合定位细节和深层语义信息。
\end{AnswerBox}
\else
\PracticeSpace{2.4cm}
\fi
\Question{22.}{U-Net为什么使用跳跃连接？}
\ifshowsolutions
\begin{AnswerBox}
\par\noindent\textbf{编码器经过多次池化后：}
\par\noindent\textbullet\quad 语义信息增强；
\par\noindent\textbullet\quad 空间分辨率降低；
\par 目标边界和定位细节丢失。
\par\noindent\textbf{编码器浅层特征保留了丰富的：}
\par\noindent\textbullet\quad 边缘；
\par\noindent\textbullet\quad 纹理；
\par 空间位置信息。
\par 跳跃连接将浅层定位信息与解码器的高级语义特征融合，从而提高分割边界和小目标的恢复能力。
\end{AnswerBox}
\else
\PracticeSpace{2.4cm}
\fi
\section{二、解答题专项}
\Question{解答题1}{卷积输出尺寸和参数量}
\par\noindent\textbf{输入特征图为：}
\par 32×32×3
\par\noindent\textbf{卷积层参数为：}
\par\noindent\textbullet\quad 卷积核大小：5×5；
\par\noindent\textbullet\quad 输出通道数：16；
\par\noindent\textbullet\quad 步长：1；
\par\noindent\textbullet\quad 填充：2；
\par 包含偏置。
\par 求输出特征图尺寸和参数量。
\ifshowsolutions
\begin{AnswerBox}
\par\noindent\textbf{输出高度：}
\begin{center}\small\ttfamily H\_out = (32 + 2×2 - 5)/1 + 1 = 32\end{center}
\par\noindent\textbf{输出宽度同样为32，输出通道数为16：}
\par 32×32×16
\par\noindent\textbf{每个卷积核参数量为：}
\par 5×5×3+1=76
\par\noindent\textbf{共有16个卷积核：}
\par 76×16=1216
\par\noindent\textbf{因此参数量为：}
\par 1216
\end{AnswerBox}
\else
\PracticeSpace{4.2cm}
\fi
\Question{解答题2}{多层CNN尺寸与参数}
\par\noindent\textbf{网络结构如下：}
\par 输入：64 × 64 × 3
\par\noindent\textbf{Conv1：}
\par 3 × 3卷积，输出通道16
\par stride=1，padding=1
\par\noindent\textbf{MaxPool：}
\par 2 × 2，stride=2
\par\noindent\textbf{Conv2：}
\par 5 × 5卷积，输出通道32
\par stride=1，padding=0
\par 求各层输出尺寸以及两个卷积层的参数量，均包含偏置。
\ifshowsolutions
\begin{AnswerBox}
\par Conv1
\par\noindent\textbf{输出尺寸：}
\begin{center}\small\ttfamily H\_out = (64 + 2 - 3)/1 + 1 = 64\end{center}
\par\noindent\textbf{所以：}
\par 64×64×16
\par\noindent\textbf{参数量：}
\begin{center}\small\ttfamily (3×3×3+1)×16 = 28×16 = 448\end{center}
\par MaxPool
\par\noindent\textbf{空间尺寸减半，通道数不变：}
\par 32×32×16
\par Conv2
\par\noindent\textbf{输出空间尺寸：}
\begin{center}\small\ttfamily H\_out = (32 - 5)/1 + 1 = 28\end{center}
\par\noindent\textbf{所以：}
\par 28×28×32
\par\noindent\textbf{参数量：}
\begin{center}\small\ttfamily (5×5×16+1)×32 =401×32=12832\end{center}
\par\noindent\textbf{因此：}
\par Conv2参数量=12832
\end{AnswerBox}
\else
\PracticeSpace{4.2cm}
\fi
\Question{解答题3}{CNN加全连接层}
\par\noindent\textbf{网络结构为：}
\par 输入：28 × 28 × 1
\par\noindent\textbf{Conv：}
\par 5 × 5卷积，输出通道6
\par stride=1，padding=0
\par\noindent\textbf{MaxPool：}
\par 2 × 2，stride=2
\par Flatten
\par\noindent\textbf{全连接层：}
\par 输出10个神经元
\par 计算各层输出尺寸和总参数量。
\ifshowsolutions
\begin{AnswerBox}
\par 卷积层输出
\par 28−5+1=24
\par\noindent\textbf{因此：}
\par 24×24×6
\par\noindent\textbf{卷积参数量：}
\par (5×5×1+1)×6=26×6=156
\par 池化层输出
\par 12×12×6
\par Flatten
\par 12×12×6=864
\par 全连接层参数量
\par 864×10+10=8650
\par 总参数量
\par 156+8650=8806
\par\noindent\textbf{最终：}
\par 总参数量=8806
\end{AnswerBox}
\else
\PracticeSpace{4.2cm}
\fi
\Question{解答题4}{判断卷积参数是否计算偏置}
\par 某卷积层输入通道为64，输出通道为128，卷积核大小为 3×3。
\par\noindent\textbf{分别计算：}
\par\noindent\textbullet\quad 包含偏置时的参数量；
\par 不包含偏置时的参数量。
\ifshowsolutions
\begin{AnswerBox}
\par\noindent\textbf{卷积权重数量：}
\par 3×3×64×128=73728
\par\noindent\textbf{包含偏置时，每个输出通道有一个偏置：}
\par 73728+128=73856
\par\noindent\textbf{因此：}
\par 无偏置
\par 含偏置
\begin{center}\small\ttfamily 无偏置 = 73728；含偏置 = 73856\end{center}
\end{AnswerBox}
\else
\PracticeSpace{4.2cm}
\fi
\Question{解答题5}{残差块尺寸匹配}
\par\noindent\textbf{残差块输入为：}
\begin{center}\small\ttfamily x ∈ R\textasciicircum{}(N×64×56×56)\end{center}
\par\noindent\textbf{主分支第一层为：}
\par 3 × 3卷积
\par 输入通道64
\par 输出通道128
\begin{center}\small\ttfamily stride=2 padding=1\end{center}
\par\noindent\textbf{主分支最终保持128个通道。回答：}
\par\noindent\textbullet\quad 主分支输出尺寸；
\par\noindent\textbullet\quad 原输入能否直接与主分支相加；
\par 捷径分支应如何处理。
\ifshowsolutions
\begin{AnswerBox}
\par\noindent\textbf{主分支空间尺寸：}
\begin{center}\small\ttfamily floor((56 + 2 - 3)/2) + 1 = 28\end{center}
\par\noindent\textbf{因此主分支输出：}
\par N×128×28×28
\par\noindent\textbf{原输入为：}
\par N×64×56×56
\par 通道数和空间尺寸均不同，不能直接相加。
\par\noindent\textbf{捷径分支应使用：}
\par\noindent\textbullet\quad 1×1 卷积；
\par\noindent\textbullet\quad 输出通道数128；
\par 步长2。
\par\noindent\textbf{得到：}
\par N×128×28×28
\par 之后再与主分支相加。
\end{AnswerBox}
\else
\PracticeSpace{4.2cm}
\fi
\Question{解答题6}{连续小卷积核的感受野}
\par 连续使用两个步长为1的 3×3 卷积，不使用池化。
\par 求第二层输出神经元相对于原图的感受野大小。
\ifshowsolutions
\begin{AnswerBox}
\par\noindent\textbf{第一层 3×3 卷积的感受野为：}
\par 3×3
\par 第二层每移动一个位置，对应第一层特征移动一个位置。第二层再使用 3×3 卷积，感受野每个方向增加2：
\par 3+2=5
\par\noindent\textbf{因此两个连续的 3×3 卷积对应：}
\par 5×5
\par 感受野。
\par 这也是VGG使用多个小卷积核替代大卷积核的原因之一。
\end{AnswerBox}
\else
\PracticeSpace{4.2cm}
\fi
\Question{解答题7}{Anchor数量计算}
\par\noindent\textbf{某特征图尺寸为：}
\par 40×30
\par\noindent\textbf{在特征图的每个位置设置：}
\par\noindent\textbullet\quad 3种尺度；
\par 3种长宽比。
\par 求Anchor总数。
\ifshowsolutions
\begin{AnswerBox}
\par\noindent\textbf{每个位置Anchor数量：}
\par 3×3=9
\par\noindent\textbf{特征图位置总数：}
\par 40×30=1200
\par\noindent\textbf{因此Anchor总数：}
\par 1200×9=10800
\par\noindent\textbf{所以：}
\par 10800
\end{AnswerBox}
\else
\PracticeSpace{4.2cm}
\fi
\Question{解答题8}{IoU计算}
\par\noindent\textbf{预测框左上角、右下角坐标为：}
\begin{center}\small\ttfamily B p =(1,1,5,5)\end{center}
\par\noindent\textbf{真实框为：}
\begin{center}\small\ttfamily B g =(3,2,7,6)\end{center}
\par 计算IoU。
\ifshowsolutions
\begin{AnswerBox}
\par\noindent\textbf{预测框宽高：}
\par 5−1=4,5−1=4
\par\noindent\textbf{面积：}
\begin{center}\small\ttfamily S\_p =16\end{center}
\par\noindent\textbf{真实框宽高：}
\par 7−3=4,6−2=4
\par\noindent\textbf{面积：}
\begin{center}\small\ttfamily S\_g =16\end{center}
\par\noindent\textbf{交集区域左上角为：}
\par (max(1,3),max(1,2))=(3,2)
\par\noindent\textbf{右下角为：}
\par (min(5,7),min(5,6))=(5,5)
\par\noindent\textbf{交集宽高：}
\par 5−3=2,5−2=3
\par\noindent\textbf{交集面积：}
\begin{center}\small\ttfamily S\_intersection =2×3=6\end{center}
\par\noindent\textbf{并集面积：}
\begin{center}\small\ttfamily S\_union =16+16-6=26\end{center}
\par\noindent\textbf{因此：}
\begin{center}\small\ttfamily IoU = 6/26 = 3/13 ≈ 0.231\end{center}
\par\noindent\textbf{所以：}
\par IoU≈0.231
\end{AnswerBox}
\else
\PracticeSpace{4.2cm}
\fi
\Question{解答题9}{NMS过程}
\par\noindent\textbf{同一类别有4个预测框：}
\par 框\quad 置信度
\begin{CodeBlock}
A    0.95
\end{CodeBlock}
\begin{CodeBlock}
B    0.88
\end{CodeBlock}
\begin{CodeBlock}
C    0.80
\end{CodeBlock}
\begin{CodeBlock}
D    0.60
\end{CodeBlock}
\par\noindent\textbf{已知：}
\begin{center}\small\ttfamily IoU(A,B)=0.70 IoU(A,C)=0.20 IoU(A,D)=0.10 IoU(C,D)=0.65\end{center}
\par NMS阈值为0.5，求最终保留的框。
\ifshowsolutions
\begin{AnswerBox}
\par 首先选择置信度最高的A并保留。
\par\noindent\textbf{B与A的IoU为：}
\par 0.70>0.5
\par 所以删除B。
\par\noindent\textbf{C与A的IoU为：}
\par 0.20<0.5
\par 保留候选资格。
\par\noindent\textbf{D与A的IoU为：}
\par 0.10<0.5
\par 也暂时保留。
\par 接下来选择剩余框中置信度最高的C并保留。
\par\noindent\textbf{D与C的IoU为：}
\par 0.65>0.5
\par 因此删除D。
\par\noindent\textbf{最终保留：}
\par A、C
\end{AnswerBox}
\else
\PracticeSpace{4.2cm}
\fi
\Question{解答题10}{ROI Pooling计算}
\par\noindent\textbf{一个ROI在特征图上对应：}
\[
X = \begin{bmatrix}
1 & 5 & 7 & 2 \\
4 & 6 & 3 & 5 \\
2 & 8 & 9 & 1 \\
3 & 2 & 4 & 6
\end{bmatrix}
\]
\par 将其ROI Pooling为 2×2，采用最大池化，求输出。
\ifshowsolutions
\begin{AnswerBox}
\par 将区域划分为4个 2×2 子区域。
\par\noindent\textbf{左上区域：}
\begin{center}\small\ttfamily [[1, 5], [4, 6]]\end{center}
\par 最大值为6。
\par\noindent\textbf{右上区域：}
\begin{center}\small\ttfamily [[2, 8], [3, 2]]\end{center}
\par 最大值为8。
\par\noindent\textbf{左下区域：}
\begin{center}\small\ttfamily [[7, 2], [3, 5]]\end{center}
\par 最大值为7。
\par\noindent\textbf{右下区域：}
\begin{center}\small\ttfamily [[9, 1], [4, 6]]\end{center}
\par 最大值为9。
\par\noindent\textbf{所以输出为：}
\begin{center}\small\ttfamily [[6, 7], [8, 9]]\end{center}
\end{AnswerBox}
\else
\PracticeSpace{4.2cm}
\fi
\Question{解答题11}{U-Net尺寸与通道计算}
\par\noindent\textbf{简化U-Net中：}
\par 输入：3 × 256 × 256
\par DoubleConv：3 → 64
\par MaxPool：2 × 2，stride=2
\par DoubleConv：64 → 128
\par 转置卷积：128 → 64，空间尺寸扩大2倍
\par 与编码器第一层特征拼接
\par DoubleConv：128 → 64
\par 1 × 1卷积：64 → 4
\par 求各主要阶段特征尺寸。
\ifshowsolutions
\begin{AnswerBox}
\par\noindent\textbf{输入：}
\par 3×256×256
\par\noindent\textbf{第一层DoubleConv保持空间尺寸：}
\par 64×256×256
\par\noindent\textbf{MaxPool后：}
\par 64×128×128
\par\noindent\textbf{第二层DoubleConv：}
\par 128×128×128
\par\noindent\textbf{转置卷积上采样并降通道：}
\par 64×256×256
\par\noindent\textbf{与编码器第一层的：}
\par 64×256×256
\par\noindent\textbf{沿通道维拼接：}
\par (64+64)×256×256
\par\noindent\textbf{即：}
\par 128×256×256
\par\noindent\textbf{解码器DoubleConv后：}
\par 64×256×256
\par\noindent\textbf{最后 1×1 卷积输出：}
\par 4×256×256
\par 表示每个像素对应4个类别分数。
\end{AnswerBox}
\else
\PracticeSpace{4.2cm}
\fi
\Question{解答题12}{训练代码顺序分析}
\par\noindent\textbf{下面训练代码存在错误：}
\begin{CodeBlock}
for images, labels in train_loader:
    outputs = model(images)
    loss = criterion(outputs, labels)
    optimizer.step()
    loss.backward()
    optimizer.zero_grad()
\end{CodeBlock}
\par 指出错误并写出正确顺序。
\ifshowsolutions
\begin{AnswerBox}
\par\noindent\textbf{原代码的问题是：}
\par\noindent\textbullet\quad 在反向传播前调用了 optimizer.step()；
\par\noindent\textbullet\quad 梯度还没有计算，无法正确更新参数；
\par 清空梯度的位置不合理。
\par\noindent\textbf{正确顺序为：}
\begin{CodeBlock}
for images, labels in train_loader:
    optimizer.zero_grad()

    outputs = model(images)
    loss = criterion(outputs, labels)

    loss.backward()
    optimizer.step()
\end{CodeBlock}
\par\noindent\textbf{对应过程：}
\par 梯度清零
\par → 前向传播
\par → 计算损失
\par → 反向传播
\par → 更新参数
\end{AnswerBox}
\else
\PracticeSpace{4.2cm}
\fi
\section{三、综合性问答题}
\Question{1.}{为什么传统方法需要人工设计特征，而CNN可以自动学习特征？}
\ifshowsolutions
\begin{AnswerBox}
\par 传统视觉方法通常需要人工设计边缘、纹理、角点、SIFT或HOG等特征，再将特征输入分类器。
\par\noindent\textbf{CNN通过多层卷积从数据中自动学习特征：}
\par\noindent\textbullet\quad 浅层学习边缘和纹理；
\par\noindent\textbullet\quad 中层学习局部形状和部件；
\par 深层学习目标语义。
\par 训练过程中，特征提取器和分类器根据同一损失函数共同优化，实现端到端学习。
\end{AnswerBox}
\else
\PracticeSpace{2.4cm}
\fi
\Question{2.}{为什么图像分类网络不能直接完成精确图像分割？}
\ifshowsolutions
\begin{AnswerBox}
\par 普通分类网络经过多次池化和下采样后，空间分辨率大幅降低，最终通常只输出整幅图像的类别。
\par\noindent\textbf{图像分割需要恢复到原图大小，并对每个像素分类，因此需要：}
\par\noindent\textbullet\quad 保留空间信息；
\par\noindent\textbullet\quad 使用上采样恢复尺寸；
\par\noindent\textbullet\quad 融合浅层定位特征；
\par 输出与输入空间对应的特征图。
\par U-Net正是通过编码器、解码器和跳跃连接实现这一过程。
\end{AnswerBox}
\else
\PracticeSpace{2.4cm}
\fi
\chapter{Faster R-CNN专项题库}
\ChapterIntro{围绕Anchor、RPN、Proposal、ROI Pooling及其计算过程进行集中训练。}
\section{一、Anchor 专项问答题}
\Question{1.}{什么是 Anchor？}
\ifshowsolutions
\begin{AnswerBox}
\par Anchor是在特征图的每一个空间位置预先设置的一组参考框。
\par\noindent\textbf{这些参考框通常具有不同的：}
\par\noindent\textbullet\quad 尺寸；
\par 长宽比。
\par\noindent\textbf{RPN以Anchor为基础，预测：}
\par\noindent\textbullet\quad 该Anchor属于前景还是背景；
\par 该Anchor应如何平移和缩放，才能更接近真实目标框。
\par Anchor不是最终检测框，而是边界框回归的初始参考。
\end{AnswerBox}
\else
\PracticeSpace{2.4cm}
\fi
\Question{2.}{为什么需要设置多个尺度和长宽比的Anchor？}
\ifshowsolutions
\begin{AnswerBox}
\par\noindent\textbf{图像中的目标可能具有不同的大小和形状，例如：}
\par\noindent\textbullet\quad 行人通常较高、较窄；
\par\noindent\textbullet\quad 汽车通常较宽；
\par\noindent\textbullet\quad 远处目标较小；
\par 近处目标较大。
\par 在同一个特征图位置设置不同尺度和长宽比的Anchor，可以使网络覆盖更多类型的目标框，提高候选框与真实目标的初始重合程度。
\end{AnswerBox}
\else
\PracticeSpace{2.4cm}
\fi
\Question{3.}{Anchor是设置在原图上，还是设置在特征图上？}
\ifshowsolutions
\begin{AnswerBox}
\par Anchor在计算上对应特征图的每个位置，但每个Anchor表示的是原始图像坐标系中的一个候选框。
\par 例如，主干网络的总步长为16，那么特征图上相邻两个位置，在原图中大约相距16个像素。
\par\noindent\textbf{因此：}
\par\noindent\textbullet\quad Anchor的中心由特征图位置决定；
\par\noindent\textbullet\quad Anchor的实际尺寸通常按照原图尺度定义；
\par 网络预测是在特征图上完成的。
\end{AnswerBox}
\else
\PracticeSpace{2.4cm}
\fi
\Question{4.}{Anchor的数量如何计算？}
\ifshowsolutions
\begin{AnswerBox}
\par\noindent\textbf{若特征图大小为：}
\par H×W
\par\noindent\textbf{每个位置设置：}
\par\noindent\textbullet\quad S 种尺度；
\par R 种长宽比。
\par\noindent\textbf{则每个位置的Anchor数为：}
\par k=S×R
\par\noindent\textbf{总Anchor数为：}
\par H×W×k
\end{AnswerBox}
\else
\PracticeSpace{2.4cm}
\fi
\Question{5.}{什么是正Anchor、负Anchor和忽略Anchor？}
\ifshowsolutions
\begin{AnswerBox}
\par 通常根据Anchor与真实框的IoU划分。
\par 正Anchor
\par\noindent\textbf{满足以下情况之一：}
\par\noindent\textbullet\quad 与某个真实框的IoU高于正样本阈值；
\par 是与某个真实框IoU最大的Anchor。
\par\noindent\textbf{正Anchor用于学习：}
\par\noindent\textbullet\quad 前景分类；
\par 边界框回归。
\par 负Anchor
\par 与所有真实框的IoU都较低，视为背景，只参与前景、背景分类。
\par 忽略Anchor
\par IoU处于正负阈值之间，不参与当前损失计算，避免模糊样本干扰训练。
\par 具体阈值可能因实现而不同。
\end{AnswerBox}
\else
\PracticeSpace{2.4cm}
\fi
\Question{6.}{为什么与真实框IoU最大的Anchor通常要被设为正样本？}
\ifshowsolutions
\begin{AnswerBox}
\par 某些真实目标可能没有任何Anchor达到规定的正样本IoU阈值。
\par 如果只按照固定阈值划分，该目标可能完全没有正样本，网络就无法学习检测它。
\par 因此，将与每个真实框IoU最大的Anchor设为正样本，可以保证每个真实目标至少对应一个正Anchor。
\end{AnswerBox}
\else
\PracticeSpace{2.4cm}
\fi
\Question{7.}{RPN对每个Anchor输出什么？}
\ifshowsolutions
\begin{AnswerBox}
\par\noindent\textbf{RPN对每个Anchor通常输出两类结果：}
\par\noindent\textbullet\quad 目标性分类结果：判断Anchor是前景还是背景；
\par 边界框回归结果：预测4个坐标偏移量。
\par\noindent\textbf{边界框回归通常包括：}
\begin{center}\small\ttfamily t\_x ,t\_y ,t\_w ,t\_h\end{center}
\par\noindent\textbf{分别控制：}
\par\noindent\textbullet\quad 中心横坐标移动；
\par\noindent\textbullet\quad 中心纵坐标移动；
\par\noindent\textbullet\quad 宽度缩放；
\par 高度缩放。
\end{AnswerBox}
\else
\PracticeSpace{2.4cm}
\fi
\Question{8.}{如果每个位置有 k 个Anchor，RPN分类分支和回归分支分别输出多少通道？}
\ifshowsolutions
\begin{AnswerBox}
\par 若二分类采用每个Anchor两个分数，即前景和背景，则分类分支输出：
\par 2k
\par 个通道。
\par\noindent\textbf{每个Anchor需要预测4个回归量，所以回归分支输出：}
\par 4k
\par 个通道。
\par\noindent\textbf{例如每个位置有9个Anchor：}
\par 分类输出通道=2×9=18
\par 回归输出通道=4×9=36
\par 有些实现使用每个Anchor一个前景分数，此时分类输出为 k 个通道。
\end{AnswerBox}
\else
\PracticeSpace{2.4cm}
\fi
\Question{9.}{Anchor与Proposal有什么区别？}
\ifshowsolutions
\begin{AnswerBox}
\par\noindent\textbf{Anchor：}
\par\noindent\textbullet\quad 是预先设置的参考框；
\par\noindent\textbullet\quad 数量很多；
\par\noindent\textbullet\quad 位置和形状尚未根据图像内容调整；
\par 是RPN预测的基础。
\par\noindent\textbf{Proposal：}
\par\noindent\textbullet\quad 是Anchor经过边界框回归调整后的候选框；
\par\noindent\textbullet\quad 通常还经过越界裁剪、低分框删除和NMS筛选；
\par\noindent\textbullet\quad 更接近真实目标；
\par 被送入ROI Pooling和后续检测头。
\par\noindent\textbf{流程为：}
\par Anchor
\par → RPN分类与回归
\par → 调整后的候选框
\par → NMS
\par → Proposal
\end{AnswerBox}
\else
\PracticeSpace{2.4cm}
\fi
\Question{10.}{为什么RPN之后还要进行NMS？}
\ifshowsolutions
\begin{AnswerBox}
\par 同一个目标附近通常会有许多相互重叠的Anchor，它们经过回归后会产生大量相似的候选框。
\par NMS保留得分最高的候选框，并删除与其高度重叠的其他框，从而：
\par\noindent\textbullet\quad 减少重复Proposal；
\par\noindent\textbullet\quad 降低后续ROI处理的计算量；
\par 保留更有代表性的候选区域。
\end{AnswerBox}
\else
\PracticeSpace{2.4cm}
\fi
\Question{11.}{Anchor过多有什么优点和缺点？}
\ifshowsolutions
\begin{AnswerBox}
\par\noindent\textbf{优点：}
\par\noindent\textbullet\quad 可以覆盖更多尺度和形状的目标；
\par\noindent\textbullet\quad 提高目标与某个Anchor获得较高IoU的概率；
\par 有利于检测尺度差异较大的目标。
\par\noindent\textbf{缺点：}
\par\noindent\textbullet\quad 增加计算量和存储量；
\par\noindent\textbullet\quad 产生大量背景负样本；
\par\noindent\textbullet\quad 正负样本更加不平衡；
\par 产生更多重复候选框。
\par 因此Anchor数量需要在覆盖能力和计算成本之间平衡。
\end{AnswerBox}
\else
\PracticeSpace{2.4cm}
\fi
\Question{12.}{为什么Anchor机制会造成正负样本不平衡？}
\ifshowsolutions
\begin{AnswerBox}
\par 特征图上会产生大量Anchor，但图像中的真实目标数量通常较少。
\par\noindent\textbf{因此：}
\par\noindent\textbullet\quad 只有少数Anchor与真实框高度重合，成为正样本；
\par 大部分Anchor都属于背景，成为负样本。
\par 训练时通常不会使用全部Anchor，而是按照一定比例采样正负样本，例如尽量保持一定的正负样本比例。
\end{AnswerBox}
\else
\PracticeSpace{2.4cm}
\fi
\section{二、ROI Pooling 专项问答题}
\Question{13.}{什么是ROI？}
\ifshowsolutions
\begin{AnswerBox}
\par ROI是Region of Interest，即感兴趣区域。
\par 在Faster R-CNN中，ROI通常指RPN生成的Proposal映射到共享特征图后所对应的区域。
\par\noindent\textbf{每个ROI可能具有不同的：}
\par\noindent\textbullet\quad 高度；
\par\noindent\textbullet\quad 宽度；
\par 位置。
\end{AnswerBox}
\else
\PracticeSpace{2.4cm}
\fi
\Question{14.}{为什么需要ROI Pooling？}
\ifshowsolutions
\begin{AnswerBox}
\par RPN产生的候选区域大小不同，但后续全连接层要求输入特征长度固定。
\par ROI Pooling可以将任意大小的候选区域转换为固定大小的特征图，例如：
\par 7×7
\par 因此不同尺寸的候选框都可以送入同一个分类和边界框回归网络。
\end{AnswerBox}
\else
\PracticeSpace{2.4cm}
\fi
\Question{15.}{简述ROI Pooling的基本步骤。}
\ifshowsolutions
\begin{AnswerBox}
\par 以输出 H
\begin{center}\small\ttfamily H\_o × W\_o\end{center}
\par\noindent\textbf{为例：}
\par\noindent\textbullet\quad 将原图中的Proposal映射到共享特征图；
\par\noindent\textbullet\quad 对ROI边界坐标进行量化取整；
\par 将ROI划分为 H
\begin{center}\small\ttfamily H\_o × W\_o\end{center}
\par\noindent\textbullet\quad 个子区域；
\par\noindent\textbullet\quad 在每个子区域内进行最大池化；
\par 得到固定大小的 H
\begin{center}\small\ttfamily H\_o × W\_o\end{center}
\par 特征图。
\par\noindent\textbf{若输入特征图有 C 个通道，最终输出为：}
\begin{center}\small\ttfamily C × H\_o × W\_o\end{center}
\end{AnswerBox}
\else
\PracticeSpace{2.4cm}
\fi
\Question{16.}{ROI Pooling是否把整个ROI缩放成固定大小？}
\ifshowsolutions
\begin{AnswerBox}
\par 从效果上看，它把不同大小的ROI转换为固定大小特征，但实现时不是简单地直接缩放图像。
\par\noindent\textbf{ROI Pooling通常是：}
\par\noindent\textbullet\quad 将ROI划分为固定数量的网格；
\par\noindent\textbullet\quad 每个网格内进行最大池化；
\par 每个网格产生一个输出值。
\par 因此，它本质上是分块最大池化。
\end{AnswerBox}
\else
\PracticeSpace{2.4cm}
\fi
\Question{17.}{ROI Pooling为什么使用共享特征图？}
\ifshowsolutions
\begin{AnswerBox}
\par 如果每个候选框都单独通过CNN，会产生大量重复卷积计算，这是R-CNN速度慢的主要原因之一。
\par\noindent\textbf{Fast R-CNN和Faster R-CNN采用：}
\par\noindent\textbullet\quad 整幅图像只经过一次主干网络；
\par 所有ROI都从同一个共享特征图中提取特征。
\par 这样可以显著减少重复计算，提高检测速度。
\end{AnswerBox}
\else
\PracticeSpace{2.4cm}
\fi
\Question{18.}{ROI Pooling的输出尺寸由什么决定？}
\ifshowsolutions
\begin{AnswerBox}
\par ROI Pooling的输出空间尺寸由预先指定的池化目标大小决定，例如：
\par 7×7
\par 而不是由输入ROI本身的大小决定。
\par 如果共享特征图通道数为 C，则无论输入ROI大小如何，输出始终为：
\par C×7×7
\end{AnswerBox}
\else
\PracticeSpace{2.4cm}
\fi
\Question{19.}{ROI Pooling会改变通道数吗？}
\ifshowsolutions
\begin{AnswerBox}
\par 通常不会。
\par ROI Pooling只对每个通道中的空间区域分别进行池化，因此：
\par\noindent\textbullet\quad 输入有 C 个通道；
\par 输出仍有 C 个通道。
\par 它改变的是空间尺寸，而不是通道数。
\end{AnswerBox}
\else
\PracticeSpace{2.4cm}
\fi
\Question{20.}{ROI Pooling的主要缺点是什么？}
\ifshowsolutions
\begin{AnswerBox}
\par\noindent\textbf{ROI Pooling需要对：}
\par\noindent\textbullet\quad ROI边界坐标；
\par 子区域边界坐标
\par 进行取整量化。
\par 这种取整会导致ROI与原始目标之间出现空间偏移，产生量化误差。
\par 对于普通目标检测问题影响可能较小，但对于需要精确像素定位的任务，例如实例分割，影响会更加明显。
\end{AnswerBox}
\else
\PracticeSpace{2.4cm}
\fi
\Question{21.}{ROI Pooling和ROI Align有什么区别？}
\ifshowsolutions
\begin{AnswerBox}
\par ROI Pooling
\par\noindent\textbullet\quad 对ROI边界取整；
\par\noindent\textbullet\quad 对分块边界取整；
\par\noindent\textbullet\quad 每个区域进行最大池化；
\par 存在明显量化误差。
\par ROI Align
\par\noindent\textbullet\quad 不进行粗糙取整；
\par\noindent\textbullet\quad 在浮点坐标位置进行采样；
\par\noindent\textbullet\quad 使用双线性插值计算特征值；
\par 空间对齐更加准确。
\par ROI Align常用于Mask R-CNN等对空间定位精度要求较高的网络。
\end{AnswerBox}
\else
\PracticeSpace{2.4cm}
\fi
\Question{22.}{ROI Pooling与普通最大池化有什么区别？}
\ifshowsolutions
\begin{AnswerBox}
\par\noindent\textbf{普通最大池化：}
\par\noindent\textbullet\quad 对整张特征图使用固定大小的滑动窗口；
\par\noindent\textbullet\quad 各位置采用相同池化规则；
\par 输出大小由窗口、步长和输入尺寸决定。
\par\noindent\textbf{ROI Pooling：}
\par\noindent\textbullet\quad 每次只处理一个ROI；
\par\noindent\textbullet\quad ROI大小可以不同；
\par\noindent\textbullet\quad 根据ROI大小动态划分子区域；
\par 所有ROI最终得到相同的输出尺寸。
\end{AnswerBox}
\else
\PracticeSpace{2.4cm}
\fi
\Question{23.}{ROI Pooling与SPP有什么联系和区别？}
\ifshowsolutions
\begin{AnswerBox}
\par\noindent\textbf{共同点：}
\par\noindent\textbullet\quad 都可以将不同大小的输入区域转换为固定长度特征；
\par 都使用空间分块和池化。
\par\noindent\textbf{区别：}
\par\noindent\textbf{SPP：}
\par\noindent\textbullet\quad 通常使用多个空间尺度；
\par\noindent\textbullet\quad 例如 1×1、2×2、4×4；
\par 将多尺度池化结果拼接。
\par\noindent\textbf{ROI Pooling：}
\par\noindent\textbullet\quad 针对每个候选区域；
\par\noindent\textbullet\quad 通常池化为单一固定尺寸，如 7×7；
\par 主要应用于Fast R-CNN和Faster R-CNN。
\end{AnswerBox}
\else
\PracticeSpace{2.4cm}
\fi
\section{三、综合问答题}
\Question{24.}{简述Faster R-CNN从输入图像到最终检测结果的完整流程。}
\ifshowsolutions
\begin{AnswerBox}
\par\noindent\textbullet\quad 输入图像经过主干卷积网络，得到共享特征图；
\par\noindent\textbullet\quad RPN在共享特征图的每个位置设置多个Anchor；
\par\noindent\textbullet\quad RPN对Anchor进行前景、背景分类和边界框回归；
\par\noindent\textbullet\quad 将调整后的候选框进行裁剪、筛选和NMS，得到Proposal；
\par\noindent\textbullet\quad 将Proposal映射到共享特征图；
\par\noindent\textbullet\quad 使用ROI Pooling将不同大小的ROI转换为固定大小特征；
\par\noindent\textbullet\quad 检测头对每个ROI进行目标类别分类；
\par\noindent\textbullet\quad 检测头进一步进行边界框回归；
\par 使用置信度筛选和NMS得到最终检测结果。
\end{AnswerBox}
\else
\PracticeSpace{2.4cm}
\fi
\Question{25.}{Anchor、RPN、Proposal和ROI Pooling之间是什么关系？}
\ifshowsolutions
\begin{AnswerBox}
\par\noindent\textbf{它们之间的关系为：}
\par 特征图每个位置生成Anchor
\par → RPN判断前景/背景并预测坐标偏移
\par → Anchor经过回归和NMS形成Proposal
\par → Proposal映射到共享特征图成为ROI
\par → ROI Pooling提取固定大小特征
\par → 分类与边界框精修
\par Anchor是初始参考框，RPN负责筛选和调整Anchor，Proposal是RPN输出的候选框，ROI Pooling负责提取固定尺寸的候选区域特征。
\end{AnswerBox}
\else
\PracticeSpace{2.4cm}
\fi
\Question{26.}{Faster R-CNN为什么被称为两阶段检测器？}
\ifshowsolutions
\begin{AnswerBox}
\par\noindent\textbf{第一阶段是RPN：}
\par\noindent\textbullet\quad 判断哪些区域可能包含目标；
\par 生成Proposal。
\par\noindent\textbf{第二阶段是检测头：}
\par\noindent\textbullet\quad 对每个Proposal进行具体类别分类；
\par 对边界框进一步精修。
\par\noindent\textbf{因此：}
\par 第一阶段：产生候选区域
\par 第二阶段：分类和精确定位
\par 所以Faster R-CNN属于两阶段目标检测器。
\end{AnswerBox}
\else
\PracticeSpace{2.4cm}
\fi
\Question{27.}{RPN的“二分类”和最终检测头的“多分类”有什么区别？}
\ifshowsolutions
\begin{AnswerBox}
\par\noindent\textbf{RPN只关心Anchor中是否存在某种目标，因此进行：}
\par 前景与背景二分类
\par 它不需要判断目标具体是人、车还是动物。
\par 最终检测头需要判断Proposal属于哪个具体类别，因此进行：
\par 背景＋多个目标类别的多分类
\par 例如检测20类目标时，检测头通常输出21类，包括20个目标类和1个背景类。
\end{AnswerBox}
\else
\PracticeSpace{2.4cm}
\fi
\Question{28.}{Faster R-CNN为什么需要进行两次边界框回归？}
\ifshowsolutions
\begin{AnswerBox}
\par\noindent\textbf{第一次回归由RPN完成：}
\par\noindent\textbullet\quad 从粗糙Anchor得到较好的Proposal；
\par 目标是找到可能存在物体的区域。
\par\noindent\textbf{第二次回归由检测头完成：}
\par\noindent\textbullet\quad 在已获得的Proposal基础上进一步精修；
\par 通常结合具体类别信息得到更准确的最终框。
\par\noindent\textbf{因此两次回归分别负责：}
\par 粗定位→精确定位
\end{AnswerBox}
\else
\PracticeSpace{2.4cm}
\fi
\section{四、计算与解答题}
\Question{解答题1}{Anchor数量计算}
\par\noindent\textbf{某共享特征图大小为：}
\par 50×40
\par 每个位置设置3种尺度和3种长宽比。
\par 求Anchor总数。
\ifshowsolutions
\begin{AnswerBox}
\par\noindent\textbf{每个位置的Anchor数量为：}
\par 3×3=9
\par\noindent\textbf{特征图位置数量为：}
\par 50×40=2000
\par\noindent\textbf{总Anchor数：}
\par 2000×9=18000
\par\noindent\textbf{因此：}
\par 18000
\end{AnswerBox}
\else
\PracticeSpace{4.2cm}
\fi
\Question{解答题2}{RPN输出通道数}
\par 特征图每个位置设置9个Anchor。
\par\noindent\textbf{假设：}
\par\noindent\textbullet\quad 分类分支对每个Anchor输出前景、背景两个分数；
\par 回归分支对每个Anchor输出4个偏移量。
\par 求两个分支输出通道数。
\ifshowsolutions
\begin{AnswerBox}
\par\noindent\textbf{分类分支：}
\par 2×9=18
\par\noindent\textbf{回归分支：}
\par 4×9=36
\par\noindent\textbf{所以：}
\par 分类分支18通道，回归分支36通道
\end{AnswerBox}
\else
\PracticeSpace{4.2cm}
\fi
\Question{解答题3}{Anchor回归计算}
\par\noindent\textbf{某Anchor为：}
\begin{center}\small\ttfamily (x\_a ,y\_a ,w\_a ,h\_a )=(40,50,20,10)\end{center}
\par\noindent\textbf{RPN预测：}
\begin{center}\small\ttfamily t\_x =0.2,t\_y =-0.1,t\_w =ln1.5,t\_h =ln2\end{center}
\par\noindent\textbf{采用：}
\begin{center}\small\ttfamily x = x\_a + t\_x w\_a，y = y\_a + t\_y h\_a，w = w\_a exp(t\_w)，h = h\_a exp(t\_h)\end{center}
\par 求调整后的候选框。
\ifshowsolutions
\begin{AnswerBox}
\par\noindent\textbf{中心横坐标：}
\par x=40+0.2×20=44
\par\noindent\textbf{中心纵坐标：}
\par y=50−0.1×10=49
\par\noindent\textbf{宽度：}
\begin{center}\small\ttfamily w = 20 exp(ln 1.5) = 20×1.5 = 30\end{center}
\par\noindent\textbf{高度：}
\begin{center}\small\ttfamily h = 10 exp(ln 2) = 10×2 = 20\end{center}
\par\noindent\textbf{所以：}
\par (x,y,w,h)=(44,49,30,20)
\end{AnswerBox}
\else
\PracticeSpace{4.2cm}
\fi
\Question{解答题4}{Anchor正负样本判断}
\par\noindent\textbf{某Anchor与三个真实框的IoU分别为：}
\par 0.20,0.72,0.35
\par\noindent\textbf{规定：}
\par\noindent\textbullet\quad 最大IoU不小于0.7为正样本；
\par\noindent\textbullet\quad 最大IoU小于0.3为负样本；
\par 其余忽略。
\par 判断该Anchor类别。
\ifshowsolutions
\begin{AnswerBox}
\par\noindent\textbf{该Anchor的最大IoU为：}
\par max(0.20,0.72,0.35)=0.72
\par\noindent\textbf{因为：}
\par 0.72≥0.7
\par 所以该Anchor为正样本。
\par 正Anchor
\end{AnswerBox}
\else
\PracticeSpace{4.2cm}
\fi
\Question{解答题5}{ROI Pooling计算}
\par\noindent\textbf{某ROI对应的特征区域为：}
\[
X = \begin{bmatrix}
1 & 5 & 2 & 4 \\
3 & 6 & 7 & 1 \\
2 & 1 & 9 & 5 \\
4 & 8 & 3 & 6
\end{bmatrix}
\]
\par 将其ROI Pooling为 2×2，采用最大池化。
\ifshowsolutions
\begin{AnswerBox}
\par\noindent\textbf{左上区域：}
\begin{center}\small\ttfamily [[1, 5], [3, 6]]\end{center}
\par 最大值为6。
\par\noindent\textbf{右上区域：}
\begin{center}\small\ttfamily [[2, 1], [4, 8]]\end{center}
\par 最大值为8。
\par\noindent\textbf{左下区域：}
\begin{center}\small\ttfamily [[2, 4], [7, 1]]\end{center}
\par 最大值为7。
\par\noindent\textbf{右下区域：}
\begin{center}\small\ttfamily [[9, 5], [3, 6]]\end{center}
\par 最大值为9。
\par\noindent\textbf{所以输出为：}
\begin{center}\small\ttfamily [[6, 7], [8, 9]]\end{center}
\end{AnswerBox}
\else
\PracticeSpace{4.2cm}
\fi
\Question{解答题6}{ROI Pooling输出尺寸}
\par 共享特征图通道数为256，对一个任意大小的ROI执行 7×7 ROI Pooling。
\par 求输出特征图尺寸以及展平后的特征数量。
\ifshowsolutions
\begin{AnswerBox}
\par\noindent\textbf{ROI Pooling不改变通道数，所以输出为：}
\par 256×7×7
\par\noindent\textbf{展平后的特征数量为：}
\par 256×7×7=12544
\par\noindent\textbf{所以：}
\par 12544
\end{AnswerBox}
\else
\PracticeSpace{4.2cm}
\fi
\backmatter
\chapter*{使用说明}
\addcontentsline{toc}{chapter}{使用说明}
本合集根据给定复习材料整理。练习版隐藏参考答案，含答案版保留完整答案与计算过程；两版题号和章节结构一致。
\end{document}
